\documentclass[a4paper,10pt]{article}
\usepackage{eumat}

\begin{document}
\begin{eulernotebook}
\eulerheading{Visualisasi dan Perhitungan Geometri dengan EMT}
\begin{eulercomment}
Catatan kuliah Geometri dengan Euler Math Toolbox (EMT).
Versi ini: (1) semua komentar sudah dalam bahasa Indonesia, (2) setiap contoh
dilengkapi uraian materi (teori), (3) semua soal Latihan dikerjakan langsung
di Euler, dan (4) setiap penyelesaian diberi penjelasan hasil.

Catatan penting: nilai yang tertulis pada bagian "Hasil yang diharapkan" dihitung
dan diperiksa terpisah (Python). Jalankan setiap perintah dengan Shift+Enter
di EMT dan bandingkan hasilnya dengan nilai tersebut.

Euler menyediakan beberapa fungsi untuk melakukan visualisasi dan perhitungan geometri, baik
secara numerik maupun analitik (seperti biasanya tentunya, menggunakan Maxima).
Fungsi-fungsi untuk visualisasi dan perhitungan geometeri tersebut disimpan di dalam file
program "geometry.e", sehingga file tersebut harus dipanggil sebelum menggunakan
fungsi-fungsi atau perintah-perintah untuk geometri.
\end{eulercomment}
\begin{eulerprompt}
>load geometry
\end{eulerprompt}
\begin{euleroutput}
  Numerical and symbolic geometry.
\end{euleroutput}
\begin{eulercomment}
\end{eulercomment}
\eulersubheading{Fungsi-fungsi Geometri}
\begin{eulercomment}
Fungsi-fungsi untuk Menggambar Objek Geometri:

\end{eulercomment}
\begin{eulerttcomment}
  defaultd:=textheight()*1.5: nilai asli untuk parameter d
  setPlotrange(x1,x2,y1,y2): menentukan rentang x dan y pada bidang koordinat
  setPlotRange(r): pusat bidang koordinat (0,0) dan batas-batas sumbu-x dan y adalah -r sd r
  plotPoint (P, "P"): menggambar titik P dan diberi label "P"
  plotSegment (A,B, "AB", d): menggambar ruas garis AB, diberi label "AB" sejauh d
  plotLine (g, "g", d): menggambar garis g diberi label "g" sejauh d
  plotCircle (c,"c",v,d): Menggambar lingkaran c dan diberi label "c"
  plotLabel (label, P, V, d): menuliskan label pada posisi P
\end{eulerttcomment}
\begin{eulercomment}

Fungsi-fungsi Geometri Analitik (numerik maupun simbolik):

\end{eulercomment}
\begin{eulerttcomment}
  turn(v, phi): memutar vektor v sejauh phi
  turnLeft(v):   memutar vektor v ke kiri
  turnRight(v):  memutar vektor v ke kanan
  normalize(v): normal vektor v
  crossProduct(v, w): hasil kali silang vektorv dan w.
  lineThrough(A, B): garis melalui A dan B, hasilnya [a,b,c] sdh. ax+by=c.
  lineWithDirection(A,v): garis melalui A searah vektor v
  getLineDirection(g): vektor arah (gradien) garis g
  getNormal(g): vektor normal (tegak lurus) garis g
  getPointOnLine(g):  titik pada garis g
  perpendicular(A, g):  garis melalui A tegak lurus garis g
  parallel (A, g):  garis melalui A sejajar garis g
  lineIntersection(g, h):  titik potong garis g dan h
  projectToLine(A, g):   proyeksi titik A pada garis g
  distance(A, B):  jarak titik A dan B
  distanceSquared(A, B):  kuadrat jarak A dan B
  quadrance(A, B): kuadrat jarak A dan B
  areaTriangle(A, B, C):  luas segitiga ABC
  computeAngle(A, B, C):   besar sudut <ABC
  angleBisector(A, B, C): garis bagi sudut <ABC
  circleWithCenter (A, r): lingkaran dengan pusat A dan jari-jari r
  getCircleCenter(c):  pusat lingkaran c
  getCircleRadius(c):  jari-jari lingkaran c
  circleThrough(A,B,C):  lingkaran melalui A, B, C
  middlePerpendicular(A, B): garis sumbu (tegak lurus di titik tengah) ruas AB
  lineCircleIntersections(g, c): titik potong garis g dan lingkran c
  circleCircleIntersections (c1, c2):  titik potong lingkaran c1 dan c2
  planeThrough(A, B, C):  bidang melalui titik A, B, C
\end{eulerttcomment}
\begin{eulercomment}

Fungsi-fungsi Khusus Untuk Geometri Simbolik:

\end{eulercomment}
\begin{eulerttcomment}
  getLineEquation (g,x,y): persamaan garis g dinyatakan dalam x dan y
  getHesseForm (g,x,y,A): bentuk Hesse garis g dinyatakan dalam x dan y dengan titik A pada
  sisi positif (kanan/atas) garis
  quad(A,B): kuadrat jarak AB
  spread(a,b,c): Spread segitiga dengan panjang sisi-sisi a,b,c, yakni sin(alpha)^2 dengan
  alpha sudut yang menghadap sisi a.
  crosslaw(a,b,c,sa): persamaan 3 quads dan 1 spread pada segitiga dengan panjang sisi a, b, c.
  triplespread(sa,sb,sc): persamaan 3 spread sa,sb,sc yang memebntuk suatu segitiga
  doublespread(sa): Spread sudut rangkap Spread 2*phi, dengan sa=sin(phi)^2 spread a.
\end{eulerttcomment}
\begin{eulercomment}

Fungsi bantu buatan sendiri (dipakai pada latihan di bawah):

\end{eulercomment}
\begin{eulerttcomment}
  jarakKeGaris(P,g): jarak titik P ke garis g=[a,b,c] (ax+by=c)
  gambarSegitiga(A,B,C): menggambar titik dan sisi-sisi segitiga ABC
\end{eulerttcomment}
\begin{eulercomment}

\end{eulercomment}
\begin{eulerprompt}
>function jarakKeGaris (P,g) := abs(g[1]*P[1]+g[2]*P[2]-g[3])/sqrt(g[1]^2+g[2]^2)
>function gambarSegitiga (A,B,C) ...
\end{eulerprompt}
\begin{eulerudf}
    plotPoint(A,"A"); plotPoint(B,"B"); plotPoint(C,"C");
    plotSegment(A,B,""); plotSegment(B,C,""); plotSegment(C,A,"");
  endfunction
\end{eulerudf}
\begin{eulercomment}
\begin{eulercomment}
\eulerheading{Uraian Materi 1: Geometri Koordinat Dasar}
\begin{eulercomment}
Titik pada bidang dinyatakan dengan pasangan terurut P(x,y), atau vektor [x,y] di EMT.

- Jarak dua titik A(x1,y1) dan B(x2,y2):
\end{eulercomment}
\begin{eulerformula}
\[
|AB| = \sqrt{(x_2-x_1)^2+(y_2-y_1)^2}.
\]
\end{eulerformula}
\begin{eulercomment}
- Persamaan garis umum: ax + by = c, dengan vektor normal (a,b). EMT menyimpan garis
\end{eulercomment}
\begin{eulerttcomment}
  sebagai vektor [a,b,c]. Garis melalui A dan B mempunyai vektor normal yang tegak lurus AB.
\end{eulerttcomment}
\begin{eulercomment}

- Dua garis tegak lurus jika hasil kali skalar vektor normalnya nol. Jika gradiennya
\end{eulercomment}
\begin{eulerttcomment}
  ada, hasil kali gradien keduanya = -1.
\end{eulerttcomment}
\begin{eulercomment}

- Jarak titik P(x0,y0) ke garis ax+by=c:
\end{eulercomment}
\begin{eulerformula}
\[
d = \frac{|a x_0 + b y_0 - c|}{\sqrt{a^2+b^2}}.
\]
\end{eulerformula}
\begin{eulerttcomment}
  Rumus inilah yang dipakai fungsi jarakKeGaris di atas.
\end{eulerttcomment}
\begin{eulercomment}
- Luas segitiga ABC dengan determinan (bertanda):
\end{eulercomment}
\begin{eulerformula}
\[
L = \frac{1}{2}\left|x_1(y_2-y_3)+x_2(y_3-y_1)+x_3(y_1-y_2)\right|.
\]
\end{eulerformula}
\begin{eulercomment}
\begin{eulercomment}
\eulerheading{Contoh 1: Luas, Lingkaran Luar, Lingkaran Dalam Segitiga}
\begin{eulercomment}
\end{eulercomment}
\eulersubheading{Uraian Materi}
\begin{eulercomment}
Garis-garis istimewa pada segitiga ABC dengan sisi a=BC, b=CA, c=AB dan s=(a+b+c)/2:

- Garis tinggi: garis dari titik sudut yang tegak lurus sisi di hadapannya. Ketiganya
\end{eulercomment}
\begin{eulerttcomment}
  berpotongan di satu titik, disebut titik tinggi (ortosentrum) H.
\end{eulerttcomment}
\begin{eulercomment}
- Garis berat: garis dari titik sudut ke titik tengah sisi di hadapannya. Ketiganya
\end{eulercomment}
\begin{eulerttcomment}
  berpotongan di titik berat G = (A+B+C)/3, yang membagi garis berat 2 : 1.
\end{eulerttcomment}
\begin{eulercomment}
- Garis sumbu: garis tegak lurus sisi di titik tengahnya. Ketiganya berpotongan di
\end{eulercomment}
\begin{eulerttcomment}
  pusat lingkaran luar O, sebab setiap titik pada garis sumbu AB berjarak sama ke A dan B.
\end{eulerttcomment}
\begin{eulercomment}
- Garis bagi sudut: ketiganya berpotongan di pusat lingkaran dalam, sebab setiap
\end{eulercomment}
\begin{eulerttcomment}
  titik pada garis bagi sudut berjarak sama ke kedua sisi yang membentuk sudut itu.
\end{eulerttcomment}
\begin{eulercomment}

Rumus-rumus yang berkaitan (L = luas segitiga):
\end{eulercomment}
\begin{eulerformula}
\[
L=\frac{1}{2}\,a\,t_a,\qquad r=\frac{L}{s},\qquad R=\frac{abc}{4L},\qquad \frac{a}{\sin\alpha}=2R.
\]
\end{eulerformula}
\begin{eulercomment}
Jari-jari lingkaran dalam r, jari-jari lingkaran luar R, dan pertidaksamaan Euler: R \textgreater{}= 2r.

\end{eulercomment}
\eulersubheading{Penyelesaian dengan EMT}
\begin{eulercomment}
Untuk menggambar objek-objek geometri, langkah pertama adalah menentukan rentang sumbu-sumbu
koordinat. Semua objek geometri akan digambar pada satu bidang koordinat, sampai didefinisikan
bidang koordinat yang baru.
\end{eulercomment}
\begin{eulerprompt}
>setPlotRange(-0.5,2.5,-0.5,2.5); // mendefinisikan bidang koordinat baru
\end{eulerprompt}
\begin{eulercomment}
Sekarang tentukan tiga titik dan gambarlah.
\end{eulercomment}
\begin{eulerprompt}
>A=[1,0]; plotPoint(A,"A"); // definisi dan gambar tiga titik
>B=[0,1]; plotPoint(B,"B");
>C=[2,2]; plotPoint(C,"C");
\end{eulerprompt}
\begin{eulercomment}
Kemudian tiga ruas garis.
\end{eulercomment}
\begin{eulerprompt}
>plotSegment(A,B,"c"); // c=AB
>plotSegment(B,C,"a"); // a=BC
>plotSegment(A,C,"b"); // b=AC
\end{eulerprompt}
\begin{eulercomment}
Fungsi-fungsi geometri mencakup fungsi untuk membuat garis dan lingkaran.
Format garis adalah [a,b,c], yang menyatakan garis dengan
persamaan ax+by=c.
\end{eulercomment}
\begin{eulerprompt}
>lineThrough(B,C) // garis yang melalui B dan C
\end{eulerprompt}
\begin{euleroutput}
  [-1,  2,  2]
\end{euleroutput}
\begin{eulercomment}
Hitung garis yang tegak lurus BC melalui A.
\end{eulercomment}
\begin{eulerprompt}
>h=perpendicular(A,lineThrough(B,C)); // garis h tegak lurus BC melalui A
\end{eulerprompt}
\begin{eulercomment}
Dan titik potongnya dengan BC.
\end{eulercomment}
\begin{eulerprompt}
>D=lineIntersection(h,lineThrough(B,C)); // D adalah titik potong h dan BC
\end{eulerprompt}
\begin{eulercomment}
Gambar itu.
\end{eulercomment}
\begin{eulerprompt}
>plotPoint(D,value=1); // koordinat D ditampilkan
>aspect(1); plotSegment(A,D): // tampilkan semua gambar hasil plot...()
\end{eulerprompt}
\begin{eulercomment}
Hitung luas ABC:

\end{eulercomment}
\begin{eulerformula}
\[
L_{\triangle ABC}= \frac{1}{2}AD.BC.
\]
\end{eulerformula}
\begin{eulerprompt}
>norm(A-D)*norm(B-C)/2 // AD=norm(A-D), BC=norm(B-C)
\end{eulerprompt}
\begin{euleroutput}
  1.5
\end{euleroutput}
\begin{eulercomment}
Bandingkan dengan rumus determinan.
\end{eulercomment}
\begin{eulerprompt}
>areaTriangle(A,B,C) // hitung luas segitiga langusng dengan fungsi
\end{eulerprompt}
\begin{euleroutput}
  1.5
\end{euleroutput}
\begin{eulercomment}
Cara lain menghitung luas segitigas ABC:
\end{eulercomment}
\begin{eulerprompt}
>distance(A,D)*distance(B,C)/2
\end{eulerprompt}
\begin{euleroutput}
  1.5
\end{euleroutput}
\begin{eulercomment}
Hasil yang diharapkan: ketiganya 1.5.

*** Latihan 1 (Contoh 1): Luas ABC dengan tinggi dari titik B dan C

Soal: Coba hitung luas ABC dengan menggunakan tinggi dari titik B maupun C.
Gambar kedua garis tinggi lainnya. Apakah hasilnya sama?

Langkah: tarik garis dari B tegak lurus sisi AC (kaki = Bk), dan dari C tegak lurus
sisi AB (kaki = Ck). Luas = 1/2 x alas x tinggi.
\end{eulercomment}
\begin{eulerprompt}
>hB=perpendicular(B,lineThrough(A,C)); Bk=lineIntersection(hB,lineThrough(A,C));
>hC=perpendicular(C,lineThrough(A,B)); Ck=lineIntersection(hC,lineThrough(A,B));
>plotPoint(Bk,"Bk"); plotPoint(Ck,"Ck"); plotSegment(B,Bk); plotSegment(C,Ck):
\end{eulerprompt}
\begin{eulercomment}
Luas dengan tinggi dari B (alas AC) dan tinggi dari C (alas AB):
\end{eulercomment}
\begin{eulerprompt}
>norm(B-Bk)*norm(A-C)/2
\end{eulerprompt}
\begin{euleroutput}
  1.5
\end{euleroutput}
\begin{eulerprompt}
>norm(C-Ck)*norm(A-B)/2
\end{eulerprompt}
\begin{euleroutput}
  1.5
\end{euleroutput}
\begin{eulercomment}
Hasil yang diharapkan: tinggi dari B = 1.34164 (=3/sqrt(5)), tinggi dari C = 2.12132
(=3/sqrt(2)); kedua luasnya 1.5.
Penjelasan: ketiga cara memberi luas yang sama (1.5). Ini wajar karena
alas x tinggi bernilai sama untuk ketiga pasangan sisi-tinggi: a.ta = b.tb = c.tc = 2L.
Tinggi yang lebih pendek berpasangan dengan alas yang lebih panjang.

Sudut di C.
\end{eulercomment}
\begin{eulerprompt}
>degprint(computeAngle(B,C,A))
\end{eulerprompt}
\begin{euleroutput}
  36°52'11.63''
\end{euleroutput}
\begin{eulercomment}
*** Latihan 2 (Contoh 1): Dua sudut lain dan jumlah ketiga sudut

Soal: Hitung besar dua sudut yang lain. Selanjutnya, jumlahkan ketiga sudut
segitiga tersebut. Apakah hasilnya 180 derajat?
\end{eulercomment}
\begin{eulerprompt}
>wA=computeAngle(B,A,C); wB=computeAngle(A,B,C); wC=computeAngle(B,C,A);
>degprint(wA), degprint(wB), degprint(wC)
\end{eulerprompt}
\begin{euleroutput}
  71°33'54.18''
  71°33'54.18''
  36°52'11.63''
\end{euleroutput}
\begin{eulerprompt}
>degprint(wA+wB+wC)
\end{eulerprompt}
\begin{euleroutput}
  180°
\end{euleroutput}
\begin{eulercomment}
Hasil yang diharapkan: A = 71°33'54.18'', B = 71°33'54.18'', C = 36°52'11.63''
dan jumlahnya 180°0'0''.
Penjelasan: jumlah sudut dalam segitiga selalu 180 derajat (geometri Euclid). Di sini sudut A
dan B sama besar karena AC = BC = sqrt(5), jadi segitiga ABC sama kaki dengan puncak C.

Sekarang lingkaran luar segitiga tersebut.
\end{eulercomment}
\begin{eulerprompt}
>c=circleThrough(A,B,C); // lingkaran luar segitiga ABC
>R=getCircleRadius(c); // jari2 lingkaran luar
>O=getCircleCenter(c); // titik pusat lingkaran c
>plotPoint(O,"O"); // gambar titik "O"
>plotCircle(c,"Lingkaran luar segitiga ABC"):
\end{eulerprompt}
\begin{eulercomment}
Tampilkan koordinat titik pusat dan jari-jari lingkaran luar.
\end{eulercomment}
\begin{eulerprompt}
>O, R
\end{eulerprompt}
\begin{euleroutput}
  [1.16667,  1.16667]
  1.17851130198
\end{euleroutput}
\begin{eulercomment}
Hasil yang diharapkan: O = [1.16667, 1.16667], R = 1.17851.
Sekarang akan digambar lingkaran dalam segitiga ABC. Titik pusat lingkaran dalam adalah
titik potong garis-garis bagi sudut.
\end{eulercomment}
\begin{eulerprompt}
>l=angleBisector(A,C,B); // garis bagi <ACB
>g=angleBisector(C,A,B); // garis bagi <CAB
>P=lineIntersection(l,g) // titik potong kedua garis bagi sudut
\end{eulerprompt}
\begin{euleroutput}
  [0.86038,  0.86038]
\end{euleroutput}
\begin{eulercomment}
Tambahkan semuanya ke gambar.
\end{eulercomment}
\begin{eulerprompt}
>color(5); plotLine(l); plotLine(g); color(1); // gambar kedua garis bagi sudut
>plotPoint(P,"P"); // gambar titik potongnya
>r=norm(P-projectToLine(P,lineThrough(A,B))) // jari-jari lingkaran dalam
\end{eulerprompt}
\begin{euleroutput}
  0.509653732104
\end{euleroutput}
\begin{eulerprompt}
>plotCircle(circleWithCenter(P,r),"Lingkaran dalam segitiga ABC"): // gambar lingkaran dalam
\end{eulerprompt}
\begin{eulercomment}
Hasil yang diharapkan: P = [0.86038, 0.86038], r = 0.509654.

\end{eulercomment}
\eulersubheading{Latihan (10 soal)}
\begin{eulercomment}
*** Latihan 1: Ketiga titik singgung lingkaran dalam dengan sisi-sisi segitiga

Titik singgung adalah kaki garis tegak lurus dari pusat lingkaran dalam P ke setiap sisi,
yaitu proyeksi P pada garis-garis sisi.
\end{eulercomment}
\begin{eulerprompt}
>sAB=lineThrough(A,B); sBC=lineThrough(B,C); sCA=lineThrough(C,A);
>T1=projectToLine(P,sAB); T2=projectToLine(P,sBC); T3=projectToLine(P,sCA);
>T1, T2, T3
\end{eulerprompt}
\begin{euleroutput}
  [0.5,  0.5]
  [0.632456,  1.31623]
  [1.31623,  0.632456]
\end{euleroutput}
\begin{eulerprompt}
>plotPoint(T1,"T1"); plotPoint(T2,"T2"); plotPoint(T3,"T3"):
\end{eulerprompt}
\begin{eulercomment}
Hasil yang diharapkan: T1 = [0.5, 0.5] (pada AB), T2 = [0.632456, 1.31623] (pada BC),
T3 = [1.31623, 0.632456] (pada CA).
Penjelasan: jari-jari lingkaran dalam tegak lurus sisi di titik singgungnya, sehingga
titik singgung = proyeksi pusat pada sisi.

*** Latihan 2: Segitiga dengan titik-titik sudut ketiga titik singgung

Gambar segitiga T1T2T3 dan selidiki jenisnya melalui panjang sisi dan sudutnya.
\end{eulercomment}
\begin{eulerprompt}
>plotSegment(T1,T2,""); plotSegment(T2,T3,""); plotSegment(T3,T1,""):
>distance(T1,T2), distance(T2,T3), distance(T3,T1)
\end{eulerprompt}
\begin{euleroutput}
  0.826905214631
  0.966999966873
  0.826905214631
\end{euleroutput}
\begin{eulerprompt}
>degprint(computeAngle(T2,T1,T3)), degprint(computeAngle(T1,T2,T3)), degprint(computeAngle(T1,T3,T2))
\end{eulerprompt}
\begin{euleroutput}
  71°33'54.18''
  54°13'2.91''
  54°13'2.91''
\end{euleroutput}
\begin{eulercomment}
Hasil yang diharapkan: sisi 0.826905, 0.967000, 0.826905; sudut 71°33'54'', 54°13'3'', 54°13'3''.
Penjelasan: segitiga singgung ini sama kaki (T1T2 = T3T1) dan lancip (semua sudut \textless{} 90°).
Secara umum sudut-sudutnya sama dengan 90° - A/2, 90° - B/2, 90° - C/2, sehingga segitiga
singgung selalu lancip. Sama kakinya berasal dari ABC yang juga sama kaki (AC = BC).

*** Latihan 3: Luas segitiga singgung
\end{eulercomment}
\begin{eulerprompt}
>LT=abs(areaTriangle(T1,T2,T3))
\end{eulerprompt}
\begin{euleroutput}
  0.324341649025
\end{euleroutput}
\begin{eulercomment}
Hasil yang diharapkan: 0.324342.
Penjelasan: luas segitiga singgung = L x r/(2R). Periksa:
\end{eulercomment}
\begin{eulerprompt}
>abs(areaTriangle(A,B,C))*r/(2*R)
\end{eulerprompt}
\begin{euleroutput}
  0.324341649025
\end{euleroutput}
\begin{eulercomment}
Hasilnya sama (0.324342).

*** Latihan 4: Garis bagi sudut ketiga melalui pusat lingkaran dalam

Garis bagi sudut di A dan C sudah dipakai untuk mendapatkan P. Sekarang garis bagi sudut di B.
\end{eulercomment}
\begin{eulerprompt}
>m=angleBisector(A,B,C); // garis bagi <ABC (di titik B)
>color(3); plotLine(m); color(1);
>m[1]*P[1]+m[2]*P[2]-m[3], jarakKeGaris(P,m)
\end{eulerprompt}
\begin{euleroutput}
  0
  0
\end{euleroutput}
\begin{eulercomment}
Hasil yang diharapkan: kedua nilai sangat dekat nol (galat pembulatan ~1e-16).
Penjelasan: P memenuhi persamaan garis m (atau berjarak 0 dari m), jadi garis bagi sudut ketiga
melalui P. Dengan demikian ketiga garis bagi sudut berpotongan di satu titik.

*** Latihan 5: Gambar jari-jari lingkaran dalam
\end{eulercomment}
\begin{eulerprompt}
>plotSegment(P,T1,""); plotSegment(P,T2,""); plotSegment(P,T3,""):
>distance(P,T1), distance(P,T2), distance(P,T3)
\end{eulerprompt}
\begin{euleroutput}
  0.509653732104
  0.509653732104
  0.509653732104
\end{euleroutput}
\begin{eulercomment}
Hasil yang diharapkan: ketiganya 0.509654, sama dengan r.

*** Latihan 6: Luas lingkaran luar dan lingkaran dalam, hubungannya dengan luas ABC
\end{eulercomment}
\begin{eulerprompt}
>LABC=abs(areaTriangle(A,B,C));
>LR=pi*R^2, Lr=pi*r^2
\end{eulerprompt}
\begin{euleroutput}
  4.36332312999
  0.81601903655
\end{euleroutput}
\begin{eulerprompt}
>LR/LABC, Lr/LABC, R/r
\end{eulerprompt}
\begin{euleroutput}
  2.90888208666
  0.544012691033
  2.31237647787
\end{euleroutput}
\begin{eulercomment}
Hasil yang diharapkan: luas lingkaran luar = 4.36332, luas lingkaran dalam = 0.816019,
LR/L = 2.90888, Lr/L = 0.544013, R/r = 2.31238.
Hubungan dengan luas segitiga (s = setengah keliling):
\end{eulercomment}
\begin{eulerprompt}
>sisiA=distance(B,C); sisiB=distance(C,A); sisiC=distance(A,B); s=(sisiA+sisiB+sisiC)/2;
>r*s, sisiA*sisiB*sisiC/(4*R), LABC
\end{eulerprompt}
\begin{euleroutput}
  1.5
  1.5
  1.5
\end{euleroutput}
\begin{eulercomment}
Hasil yang diharapkan: ketiganya 1.5.
Penjelasan: L = r.s dan L = abc/(4R). Akibatnya luas lingkaran dalam = pi.L\textasciicircum{}2/s\textasciicircum{}2 dan luas
lingkaran luar = pi.a\textasciicircum{}2.b\textasciicircum{}2.c\textasciicircum{}2/(16 L\textasciicircum{}2), yaitu hubungannya melalui L, tetapi bukan
perbandingan tetap. Hanya R/r yang selalu memenuhi R \textgreater{}= 2r (pertidaksamaan Euler); di sini
2.31238 \textgreater{}= 2 dan kesamaan terjadi hanya untuk segitiga sama sisi.

*** Latihan 7: Ketiga garis sumbu berpotongan di satu titik
\end{eulercomment}
\begin{eulerprompt}
>gA=middlePerpendicular(B,C); gB=middlePerpendicular(C,A); gC=middlePerpendicular(A,B);
>setPlotRange(-0.5,2.5,-0.5,2.5); aspect(1); gambarSegitiga(A,B,C);
>color(5); plotLine(gA); plotLine(gB); plotLine(gC); color(1);
>U=lineIntersection(gA,gB); plotPoint(U,"U"):
>U, jarakKeGaris(U,gC)
\end{eulerprompt}
\begin{euleroutput}
  [1.16667,  1.16667]
  0
\end{euleroutput}
\begin{eulercomment}
Hasil yang diharapkan: U = [1.16667, 1.16667] (sama dengan pusat lingkaran luar O), jarak U ke gC = 0 (~1e-16).
Penjelasan: titik potong dua garis sumbu juga terletak pada garis sumbu ketiga, jadi ketiganya
sepusat di U, yang sama dengan O. Alasan teorinya: U berjarak sama ke A, B, C.

*** Latihan 8: Ketiga garis bagi sudut berpotongan di satu titik
\end{eulercomment}
\begin{eulerprompt}
>bA=angleBisector(B,A,C); bB=angleBisector(A,B,C); bC=angleBisector(A,C,B);
>setPlotRange(-0.5,2.5,-0.5,2.5); aspect(1); gambarSegitiga(A,B,C);
>color(5); plotLine(bA); plotLine(bB); plotLine(bC); color(1);
>P2=lineIntersection(bA,bB); plotPoint(P2,"P"):
>P2, jarakKeGaris(P2,bC)
\end{eulerprompt}
\begin{euleroutput}
  [0.86038,  0.86038]
  0
\end{euleroutput}
\begin{eulercomment}
Hasil yang diharapkan: P2 = [0.86038, 0.86038] (sama dengan pusat lingkaran dalam P), jarak ke bC = 0.

*** Latihan 9: Ketiga garis tinggi berpotongan di satu titik
\end{eulercomment}
\begin{eulerprompt}
>hA=perpendicular(A,lineThrough(B,C)); hB=perpendicular(B,lineThrough(C,A)); hC=perpendicular(C,lineThrough(A,B));
>setPlotRange(-0.5,2.5,-0.5,2.5); aspect(1); gambarSegitiga(A,B,C);
>color(5); plotLine(hA); plotLine(hB); plotLine(hC); color(1);
>H=lineIntersection(hA,hB); plotPoint(H,"H"):
>H, jarakKeGaris(H,hC)
\end{eulerprompt}
\begin{euleroutput}
  [0.666667,  0.666667]
  0
\end{euleroutput}
\begin{eulercomment}
Hasil yang diharapkan: H = [0.666667, 0.666667], jarak ke hC = 0.

*** Latihan 10: Ketiga garis berat berpotongan di satu titik
\end{eulercomment}
\begin{eulerprompt}
>mA=lineThrough(A,(B+C)/2); mB=lineThrough(B,(C+A)/2); mC=lineThrough(C,(A+B)/2);
>setPlotRange(-0.5,2.5,-0.5,2.5); aspect(1); gambarSegitiga(A,B,C);
>color(5); plotLine(mA); plotLine(mB); plotLine(mC); color(1);
>G=lineIntersection(mA,mB); plotPoint(G,"G"):
>G, jarakKeGaris(G,mC), (A+B+C)/3
\end{eulerprompt}
\begin{euleroutput}
  [1,  1]
  0
  [1,  1]
\end{euleroutput}
\begin{eulercomment}
Hasil yang diharapkan: G = [1, 1] = (A+B+C)/3, jarak ke mC = 0.
Penjelasan (7-10): keempat pasang garis istimewa memang konkuren (Teorema Ceva). Pemeriksaan
numerik di atas menunjukkan titik potong dua garis jatuh tepat pada garis ketiga. Bukti
eksak (bukan numerik) memakai Maxima ada di Latihan Contoh 2 di bawah.

\begin{eulercomment}
\eulerheading{Contoh 2: Geometri Simbolik}
\begin{eulercomment}
\end{eulercomment}
\eulersubheading{Uraian Materi}
\begin{eulercomment}
Geometri analitik simbolik mengerjakan perhitungan yang sama secara eksak (pecahan dan akar),
bukan hampiran desimal. Keuntungannya: kesamaan (misal tiga garis konkuren) bisa dibuktikan
persis. Perintah dengan awalan "\&" dan "\textdollar{}" dijalankan oleh Maxima.

\end{eulercomment}
\eulersubheading{Penyelesaian dengan EMT}
\begin{eulercomment}
Kita bisa menghitung geometri eksak dan simbolik dengan Maxima.

File geometry.e menyediakan fungsi yang sama (dan lebih banyak) di Maxima.
Namun, sekarang kita dapat memakai perhitungan simbolik.
\end{eulercomment}
\begin{eulerprompt}
>A &= [1,0]; B &= [0,1]; C &= [2,2]; // menentukan tiga titik A, B, C
\end{eulerprompt}
\begin{eulercomment}
Fungsi-fungsi untuk garis dan lingkaran bekerja persis seperti fungsi Euler,
tetapi menghasilkan perhitungan simbolik.
\end{eulercomment}
\begin{eulerprompt}
>c &= lineThrough(B,C) // c=BC
\end{eulerprompt}
\begin{euleroutput}
  
                               [- 1, 2, 2]
  
\end{euleroutput}
\begin{eulercomment}
Kita bisa mendapatkan persamaan suatu garis dengan mudah.
\end{eulercomment}
\begin{eulerprompt}
>$getLineEquation(c,x,y), $solve(%,y) | expand // persamaan garis c
>$getLineEquation(lineThrough([x1,y1],[x2,y2]),x,y), $solve(%,y) // persamaan garis melalui(x1, y1) dan (x2, y2)
>$getLineEquation(lineThrough(A,[x1,y1]),x,y) // persamaan garis melalui A dan (x1, y1)
>h &= perpendicular(A,lineThrough(B,C)) // h melalui A tegak lurus BC
\end{eulerprompt}
\begin{euleroutput}
  
                                [2, 1, 2]
  
\end{euleroutput}
\begin{eulerprompt}
>Q &= lineIntersection(c,h) // Q titik potong garis c=BC dan h
\end{eulerprompt}
\begin{euleroutput}
  
                                   2  6
                                  [-, -]
                                   5  5
  
\end{euleroutput}
\begin{eulerprompt}
>$projectToLine(A,lineThrough(B,C)) // proyeksi A pada BC
>$distance(A,Q) // jarak AQ
>cc &= circleThrough(A,B,C); $cc // (titik pusat dan jari-jari) lingkaran melalui A, B, C
>r&=getCircleRadius(cc); $r , $float(r) // tampilkan nilai jari-jari
>$computeAngle(A,C,B) // nilai <ACB
>$solve(getLineEquation(angleBisector(A,C,B),x,y),y)[1] // persamaan garis bagi <ACB
>P &= lineIntersection(angleBisector(A,C,B),angleBisector(C,B,A)); $P // titik potong 2 garis bagi sudut
>P() // hasilnya sama dengan perhitungan sebelumnya
\end{eulerprompt}
\begin{euleroutput}
  [0.86038,  0.86038]
\end{euleroutput}
\begin{eulercomment}
Hasil yang diharapkan (bandingkan Contoh 1): c = [-1,2,2]; h = [2,1,2]; Q = [2/5,6/5];
AQ = 3/sqrt(5); pusat luar [7/6,7/6] dan R = 5/(3 sqrt(2)); sudut ACB = arccos(4/5);
garis bagi \textless{}ACB: y = x; P = [(sqrt(2)sqrt(5)+2)/6, (sqrt(2)sqrt(5)+2)/6].

*** Latihan Contoh 2: Pembuktian eksak keempat garis istimewa konkuren

Untuk menghindari kerancuan nama dengan variabel di atas, nama-nama baru memakai awalan s (simbolik).

Soal 1. Ketiga garis tinggi.
\end{eulercomment}
\begin{eulerprompt}
>shA &= perpendicular(A,lineThrough(B,C)); shB &= perpendicular(B,lineThrough(C,A)); shC &= perpendicular(C,lineThrough(A,B));
>$getLineEquation(shA,x,y), $getLineEquation(shB,x,y), $getLineEquation(shC,x,y)
>HH &= lineIntersection(shA,shB); $HH
>$getLineEquation(shC,x,y) with [x=HH[1],y=HH[2]]
\end{eulerprompt}
\begin{eulercomment}
Hasil yang diharapkan: persamaan 2x+y=2, x+2y=2, x-y=0 (bentuk penulisan Maxima boleh berbeda
urutannya); HH = [2/3,2/3]; substitusi ke garis ketiga menghasilkan kesamaan yang benar (0=0).
Penjelasan: titik potong dua garis tinggi memenuhi persamaan garis tinggi ketiga, jadi ketiganya
berpotongan di satu titik (titik tinggi) HH = (2/3, 2/3).

Soal 2. Ketiga garis bagi sudut.
\end{eulercomment}
\begin{eulerprompt}
>sbA &= angleBisector(B,A,C); sbB &= angleBisector(A,B,C); sbC &= angleBisector(A,C,B);
>PP &= lineIntersection(sbA,sbB)|radcan; $PP
>$radcan(getLineEquation(sbC,x,y) with [x=PP[1],y=PP[2]])
\end{eulerprompt}
\begin{eulercomment}
Pemeriksaan numerik sebagai cadangan (hasilnya harus sekitar 0):
\end{eulercomment}
\begin{eulerprompt}
>PP(), jarakKeGaris(PP(),sbC())
\end{eulerprompt}
\begin{euleroutput}
  [0.86038,  0.86038]
  0
\end{euleroutput}
\begin{eulercomment}
Hasil yang diharapkan: PP = [(sqrt(2)sqrt(5)+2)/6, (sqrt(2)sqrt(5)+2)/6] atau [0.86038, 0.86038];
hasil substitusi berupa kesamaan yang benar; cek numerik sekitar 0.
Penjelasan: bentuk eksaknya memuat akar karena garis bagi sudut melibatkan panjang sisi
(sqrt(2) dan sqrt(5)). Walaupun rumit, ketiga garis bagi tetap bertemu di satu titik.

Soal 3. Ketiga garis sumbu.
\end{eulercomment}
\begin{eulerprompt}
>ssA &= middlePerpendicular(B,C); ssB &= middlePerpendicular(C,A); ssC &= middlePerpendicular(A,B);
>UU &= lineIntersection(ssA,ssB); $UU
>$getLineEquation(ssC,x,y) with [x=UU[1],y=UU[2]]
\end{eulerprompt}
\begin{eulercomment}
Hasil yang diharapkan: UU = [7/6,7/6]; garis sumbu ketiga x-y=0 dipenuhi (kesamaan benar).
Penjelasan: UU = pusat lingkaran luar yang dicari pada Contoh 1 (7/6, 7/6), cocok dengan R = 5/(3 sqrt 2).

Soal 4. Ketiga garis berat.
\end{eulercomment}
\begin{eulerprompt}
>smA &= lineThrough(A,(B+C)/2); smB &= lineThrough(B,(C+A)/2); smC &= lineThrough(C,(A+B)/2);
>GG &= lineIntersection(smA,smB); $GG
>$getLineEquation(smC,x,y) with [x=GG[1],y=GG[2]]
\end{eulerprompt}
\begin{eulercomment}
Hasil yang diharapkan: GG = [1,1] = (A+B+C)/3; kesamaan benar.
Penjelasan: titik berat adalah rata-rata koordinat titik sudut. Garis berat ketiga y = x melalui (1,1).

\end{eulercomment}
\eulersubheading{Perpotongan Garis dan Lingkaran}
\begin{eulercomment}
Tentu saja, kita juga dapat memotongkan garis dengan lingkaran, dan lingkaran dengan
lingkaran.
\end{eulercomment}
\begin{eulerprompt}
>A &:= [1,0]; c=circleWithCenter(A,4);
>B &:= [1,2]; C &:= [2,1]; l=lineThrough(B,C);
>setPlotRange(5); plotCircle(c); plotLine(l);
\end{eulerprompt}
\begin{eulercomment}
Perpotongan garis dengan lingkaran menghasilkan dua titik dan banyaknya
titik potong.
\end{eulercomment}
\begin{eulerprompt}
>\{P1,P2,f\}=lineCircleIntersections(l,c);
>P1, P2, f
\end{eulerprompt}
\begin{euleroutput}
  [4.64575,  -1.64575]
  [-0.645751,  3.64575]
  2
\end{euleroutput}
\begin{eulerprompt}
>plotPoint(P1); plotPoint(P2):
\end{eulerprompt}
\begin{eulercomment}
Hal yang sama di Maxima.
\end{eulercomment}
\begin{eulerprompt}
>c &= circleWithCenter(A,4) // lingkaran dengan pusat A jari-jari 4
\end{eulerprompt}
\begin{euleroutput}
  
                                [1, 0, 4]
  
\end{euleroutput}
\begin{eulerprompt}
>l &= lineThrough(B,C) // garis l melalui B dan C
\end{eulerprompt}
\begin{euleroutput}
  
                                [1, 1, 3]
  
\end{euleroutput}
\begin{eulerprompt}
>$lineCircleIntersections(l,c) | radcan, // titik potong lingkaran c dan garis l
\end{eulerprompt}
\begin{eulercomment}
Akan ditunjukkan bahwa sudut-sudut yang menghadap busur yang sama adalah sama besar.
\end{eulercomment}
\begin{eulerprompt}
>C=A+normalize([-2,-3])*4; plotPoint(C); plotSegment(P1,C); plotSegment(P2,C);
>degprint(computeAngle(P1,C,P2))
\end{eulerprompt}
\begin{euleroutput}
  69°17'42.68''
\end{euleroutput}
\begin{eulerprompt}
>C=A+normalize([-4,-3])*4; plotPoint(C); plotSegment(P1,C); plotSegment(P2,C);
>degprint(computeAngle(P1,C,P2))
\end{eulerprompt}
\begin{euleroutput}
  69°17'42.68''
\end{euleroutput}
\begin{eulerprompt}
>insimg;
\end{eulerprompt}
\begin{eulercomment}
Hasil yang diharapkan: kedua sudut 69°17'42.68''.
Penjelasan: ini teorema sudut keliling: semua sudut keliling yang menghadap busur P1P2 sama besar,
yaitu setengah sudut pusatnya.

\end{eulercomment}
\eulersubheading{Garis Sumbu}
\begin{eulercomment}
Berikut adalah langkah-langkah menggambar garis sumbu ruas garis AB:

1. Gambar lingkaran dengan pusat A melalui B.
2. Gambar lingkaran dengan pusat B melalui A.
3. Tarik garis melallui kedua titik potong kedua lingkaran tersebut. Garis ini merupakan
garis sumbu (melalui titik tengah dan tegak lurus) AB.
\end{eulercomment}
\begin{eulerprompt}
>A=[2,2]; B=[-1,-2];
>c1=circleWithCenter(A,distance(A,B));
>c2=circleWithCenter(B,distance(A,B));
>\{P1,P2,f\}=circleCircleIntersections(c1,c2);
>l=lineThrough(P1,P2);
>setPlotRange(5); plotCircle(c1); plotCircle(c2);
>plotPoint(A); plotPoint(B); plotSegment(A,B); plotLine(l):
\end{eulerprompt}
\begin{eulercomment}
Selanjutnya, kita lakukan hal yang sama di Maxima dengan koordinat umum.
\end{eulercomment}
\begin{eulerprompt}
>A &= [a1,a2]; B &= [b1,b2];
>c1 &= circleWithCenter(A,distance(A,B));
>c2 &= circleWithCenter(B,distance(A,B));
>P &= circleCircleIntersections(c1,c2); P1 &= P[1]; P2 &= P[2];
\end{eulerprompt}
\begin{eulercomment}
Persamaan untuk titik-titik potongnya cukup rumit. Tetapi kita bisa
menyederhanakannya, jika kita selesaikan untuk y.
\end{eulercomment}
\begin{eulerprompt}
>g &= getLineEquation(lineThrough(P1,P2),x,y);
>$solve(g,y)
\end{eulerprompt}
\begin{eulercomment}
Ini memang sama dengan garis sumbu, yang dihitung dengan cara yang
sama sekali berbeda.
\end{eulercomment}
\begin{eulerprompt}
>$solve(getLineEquation(middlePerpendicular(A,B),x,y),y)
>h &=getLineEquation(lineThrough(A,B),x,y);
>$solve(h,y)
\end{eulerprompt}
\begin{eulercomment}
Perhatikan hasil kali gradien garis g dan h adalah:

\end{eulercomment}
\begin{eulerformula}
\[
\frac{-(b_1-a_1)}{(b_2-a_2)}\times \frac{(b_2-a_2)}{(b_1-a_1)} = -1.
\]
\end{eulerformula}
\begin{eulercomment}
Artinya kedua garis tegak lurus.

\begin{eulercomment}
\eulerheading{Contoh 3: Rumus Heron}
\begin{eulercomment}
\end{eulercomment}
\eulersubheading{Uraian Materi}
\begin{eulercomment}
Rumus Heron menyatakan bahwa luas segitiga dengan panjang sisi-sisi a, b dan c adalah:

\end{eulercomment}
\begin{eulerformula}
\[
L = \sqrt{s(s-a)(s-b)(s-c)}\quad \text{ dengan } s=(a+b+c)/2,
\]
\end{eulerformula}
\begin{eulercomment}
atau bisa ditulis dalam bentuk lain:

\end{eulercomment}
\begin{eulerformula}
\[
L = \frac{1}{4}\sqrt{(a+b+c)(b+c-a)(a+c-b)(a+b-c)}
\]
\end{eulerformula}
\begin{eulercomment}
Untuk membuktikan hal ini kita misalkan C(0,0), B(a,0) dan A(x,y), b=AC, c=AB. Luas segitiga
ABC adalah

\end{eulercomment}
\begin{eulerformula}
\[
L_{\triangle ABC}=\frac{1}{2}a\times y.
\]
\end{eulerformula}
\begin{eulercomment}
Nilai y didapat dengan menyelesaikan sistem persamaan:

\end{eulercomment}
\begin{eulerformula}
\[
x^2+y^2=b^2, \quad (x-a)^2+y^2=c^2.
\]
\end{eulerformula}
\begin{eulercomment}
Mengurangkan kedua persamaan memberi x = (a\textasciicircum{}2+b\textasciicircum{}2-c\textasciicircum{}2)/(2a), lalu y\textasciicircum{}2 = b\textasciicircum{}2 - x\textasciicircum{}2.

Catatan penting (perbaikan): pada versi asli bagian ini menghasilkan galat karena variabel a, b, c
masih menyimpan nilai dari contoh sebelumnya (misalnya c = [1,0,4]). Karena itu di sini semua
variabel Maxima dibersihkan dulu dengan \&remvalue().
\end{eulercomment}
\begin{eulerprompt}
>&remvalue();
>setPlotRange(-1,10,-1,8); plotPoint([0,0], "C(0,0)"); plotPoint([5.5,0], "B(a,0)");  ...
> plotPoint([7.5,6], "A(x,y)");
>plotSegment([0,0],[5.5,0], "a",25); plotSegment([5.5,0],[7.5,6],"c",15);  ...
>plotSegment([0,0],[7.5,6],"b",25);
>plotSegment([7.5,6],[7.5,0],"t=y",25):
>&assume(a>0); sol &= solve([x^2+y^2=b^2,(x-a)^2+y^2=c^2],[x,y])
\end{eulerprompt}
\begin{euleroutput}
  
                                    []
  
\end{euleroutput}
\begin{eulercomment}
Ambil solusi y.
\end{eulercomment}
\begin{eulerprompt}
>ysol &= y with sol[2][2]; $'y=sqrt(factor(ysol^2))
\end{eulerprompt}
\begin{euleroutput}
  Maxima said:
  part: invalid index of list or matrix.
   -- an error. To debug this try: debugmode(true);
  
  Error in:
  ysol &= y with sol[2][2]; $'y=sqrt(factor(ysol^2)) ...
                          ^
\end{euleroutput}
\begin{eulercomment}
Kita memperoleh rumus Heron.
\end{eulercomment}
\begin{eulerprompt}
>function H(a,b,c) &= sqrt(factor((ysol*a/2)^2)); $'H(a,b,c)=H(a,b,c)
>$'Luas=H(2,5,6) // luas segitiga dengan panjang sisi-sisi 2, 5, 6
\end{eulerprompt}
\begin{eulercomment}
Tentu saja, setiap segitiga siku-siku adalah kasus yang sudah dikenal.
\end{eulercomment}
\begin{eulerprompt}
>H(3,4,5) //luas segitiga siku-siku dengan panjang sisi 3, 4, 5
\end{eulerprompt}
\begin{euleroutput}
  Variable or function ysol not found.
  Try "trace errors" to inspect local variables after errors.
  H:
      useglobal; return a*abs(ysol)/2 
  Error in:
  H(3,4,5) //luas segitiga siku-siku dengan panjang sisi 3, 4, 5 ...
          ^
\end{euleroutput}
\begin{eulercomment}
Dan juga jelas bahwa ini adalah segitiga dengan luas maksimal
dan dua sisi 3 dan 4.
\end{eulercomment}
\begin{eulerprompt}
>aspect (1.5); plot2d(&H(3,4,x),1,7): // Kurva luas segitiga sengan panjang sisi 3, 4, x (1<= x <=7)
\end{eulerprompt}
\begin{euleroutput}
  Variable or function ysol not found.
  Error in expression: 3*abs(ysol)/2
   %ploteval:
      y0=f$(x[1],args());
  adaptiveevalone:
      s=%ploteval(g$,t;args());
  Try "trace errors" to inspect local variables after errors.
  plot2d:
      dw/n,dw/n^2,dw/n,auto;args());
\end{euleroutput}
\begin{eulercomment}
Hasil yang diharapkan: H(3,4,5) = 6. Kurva luas mencapai maksimum 6 di x = 5.
Penjelasan: luas dengan dua sisi tetap 3 dan 4 sama dengan (1/2).3.4.sin(sudut apit), maksimum
bila sudut apit = 90°, yaitu x = 5.
Kasus umum juga bekerja.
\end{eulercomment}
\begin{eulerprompt}
>$solve(diff(H(a,b,c)^2,c)=0,c)
\end{eulerprompt}
\begin{euleroutput}
  Maxima said:
  diff: second argument must be a variable; found [1,0,4]
   -- an error. To debug this try: debugmode(true);
  
  Error in:
   $solve(diff(H(a,b,c)^2,c)=0,c) ...
                                ^
\end{euleroutput}
\begin{eulercomment}
Sekarang mari kita cari himpunan semua titik di mana b+c=d untuk suatu konstanta d.
Sudah diketahui bahwa ini adalah elips.
\end{eulercomment}
\begin{eulerprompt}
>s1 &= subst(d-c,b,sol[2]); $s1
\end{eulerprompt}
\begin{euleroutput}
  Maxima said:
  part: invalid index of list or matrix.
   -- an error. To debug this try: debugmode(true);
  
  Error in:
  s1 &= subst(d-c,b,sol[2]); $s1 ...
                           ^
\end{euleroutput}
\begin{eulercomment}
Dan buatlah fungsi dari ini.
\end{eulercomment}
\begin{eulerprompt}
>function fx(a,c,d) &= rhs(s1[1]); $fx(a,c,d), function fy(a,c,d) &= rhs(s1[2]); $fy(a,c,d)
\end{eulerprompt}
\begin{eulercomment}
Sekarang kita dapat menggambar himpunan itu. Sisi b bervariasi dari 1 sampai 4.
Sudah diketahui bahwa kita mendapatkan elips.
\end{eulercomment}
\begin{eulerprompt}
>aspect(1); plot2d(&fx(3,x,5),&fy(3,x,5),xmin=1,xmax=4,square=1):
\end{eulerprompt}
\begin{eulercomment}
Kita dapat memeriksa persamaan umum elips ini, yaitu

\end{eulercomment}
\begin{eulerformula}
\[
\frac{(x-x_m)^2}{u^2}+\frac{(y-y_m)}{v^2}=1,
\]
\end{eulerformula}
\begin{eulercomment}
dengan (xm,ym) adalah pusat, dan u dan v adalah sumbu-sumbu setengah.
\end{eulercomment}
\begin{eulerprompt}
>$ratsimp((fx(a,c,d)-a/2)^2/u^2+fy(a,c,d)^2/v^2 with [u=d/2,v=sqrt(d^2-a^2)/2])
\end{eulerprompt}
\begin{eulercomment}
Kita lihat bahwa tinggi dan dengan demikian luas segitiga maksimal
untuk x=0. Jadi luas segitiga dengan a+b+c=d maksimal jika segitiga
itu sama sisi. Kita ingin menurunkan hal ini secara analitik.
\end{eulercomment}
\begin{eulerprompt}
>eqns &= [diff(H(a,b,d-(a+b))^2,a)=0,diff(H(a,b,d-(a+b))^2,b)=0]; $eqns
\end{eulerprompt}
\begin{eulercomment}
Kita mendapatkan beberapa minimum, yang berkaitan dengan segitiga dengan satu sisi 0, dan
solusi a=b=c=d/3.
\end{eulercomment}
\begin{eulerprompt}
>$solve(eqns,[a,b])
\end{eulerprompt}
\begin{eulercomment}
Ada juga metode Lagrange, memaksimumkan H(a,b,c)\textasciicircum{}2 terhadap a+b+c=d.
\end{eulercomment}
\begin{eulerprompt}
>&solve([diff(H(a,b,c)^2,a)=la,diff(H(a,b,c)^2,b)=la, ...
>   diff(H(a,b,c)^2,c)=la,a+b+c=d],[a,b,c,la])
\end{eulerprompt}
\begin{euleroutput}
  Maxima said:
  diff: second argument must be a variable; found [1,0,4]
   -- an error. To debug this try: debugmode(true);
  
  Error in:
  ... la,    diff(H(a,b,c)^2,c)=la,a+b+c=d],[a,b,c,la]) ...
                                                       ^
\end{euleroutput}
\begin{eulercomment}
Kita dapat membuat gambar dari situasi ini.
Pertama, tentukan titik-titik di Maxima.
\end{eulercomment}
\begin{eulerprompt}
>A &= at([x,y],sol[2]); $A
\end{eulerprompt}
\begin{euleroutput}
  Maxima said:
  part: invalid index of list or matrix.
   -- an error. To debug this try: debugmode(true);
  
  Error in:
  A &= at([x,y],sol[2]); $A ...
                       ^
\end{euleroutput}
\begin{eulerprompt}
>B &= [0,0]; $B, C &= [a,0]; $C
\end{eulerprompt}
\begin{eulercomment}
Kemudian tentukan rentang gambar, lalu gambar titik-titiknya.
\end{eulercomment}
\begin{eulerprompt}
>setPlotRange(0,5,-2,3); ...
>a=4; b=3; c=2; ...
>plotPoint(mxmeval("B"),"B"); plotPoint(mxmeval("C"),"C"); ...
>plotPoint(mxmeval("A"),"A"):
\end{eulerprompt}
\begin{euleroutput}
  Variable a1 not found!
  Use global variables or parameters for string evaluation.
  Error in Evaluate, superfluous characters found.
  Try "trace errors" to inspect local variables after errors.
  mxmeval:
      return evaluate(mxm(s));
  Error in:
  ... otPoint(mxmeval("C"),"C"); plotPoint(mxmeval("A"),"A"): ...
                                                       ^
\end{euleroutput}
\begin{eulercomment}
Gambar ruas-ruas garisnya.
\end{eulercomment}
\begin{eulerprompt}
>plotSegment(mxmeval("A"),mxmeval("C")); ...
>plotSegment(mxmeval("B"),mxmeval("C")); ...
>plotSegment(mxmeval("B"),mxmeval("A")):
\end{eulerprompt}
\begin{euleroutput}
  Variable a1 not found!
  Use global variables or parameters for string evaluation.
  Error in Evaluate, superfluous characters found.
  Try "trace errors" to inspect local variables after errors.
  mxmeval:
      return evaluate(mxm(s));
  Error in:
  plotSegment(mxmeval("A"),mxmeval("C")); plotSegment(mxmeval("B ...
                          ^
\end{euleroutput}
\begin{eulercomment}
Hitung garis sumbu di Maxima.
\end{eulercomment}
\begin{eulerprompt}
>h &= middlePerpendicular(A,B); g &= middlePerpendicular(B,C);
\end{eulerprompt}
\begin{eulercomment}
Dan pusat lingkaran luarnya.
\end{eulercomment}
\begin{eulerprompt}
>U &= lineIntersection(h,g);
\end{eulerprompt}
\begin{eulercomment}
Kita mendapatkan rumus untuk jari-jari lingkaran luar.
\end{eulercomment}
\begin{eulerprompt}
>&assume(a>0,b>0,c>0); $distance(U,B) | radcan
\end{eulerprompt}
\begin{eulercomment}
Mari kita tambahkan ini ke gambar.
\end{eulercomment}
\begin{eulerprompt}
>plotPoint(U()); ...
>plotCircle(circleWithCenter(mxmeval("U"),mxmeval("distance(U,C)"))):
\end{eulerprompt}
\begin{euleroutput}
  Variable a2 not found!
  Use global variables or parameters for string evaluation.
  Error in ^
  Error in expression: [a/2,(a2^2+a1^2-a*a1)/(2*a2)]
  Error in:
  plotPoint(U()); plotCircle(circleWithCenter(mxmeval("U"),mxmev ...
               ^
\end{euleroutput}
\begin{eulercomment}
Dengan geometri, kita turunkan rumus sederhana

\end{eulercomment}
\begin{eulerformula}
\[
\frac{a}{\sin(\alpha)}=2r
\]
\end{eulerformula}
\begin{eulercomment}
untuk jari-jari tersebut. Kita bisa memeriksa, apakah ini benar dengan Maxima. Maxima hanya akan
memfaktorkannya jika kita kuadratkan.
\end{eulercomment}
\begin{eulerprompt}
>$c^2/sin(computeAngle(A,B,C))^2  | factor
\end{eulerprompt}
\begin{eulercomment}
\begin{eulercomment}
\eulerheading{Contoh 4: Garis Euler dan Parabola}
\begin{eulercomment}
\end{eulercomment}
\eulersubheading{Uraian Materi}
\begin{eulercomment}
Garis Euler: pada setiap segitiga yang tidak sama sisi, titik tinggi H, titik berat G, dan
pusat lingkaran luar O terletak segaris, dan GH = 2 GO (G membagi OH dengan perbandingan 1 : 2).
Pada garis ini juga terletak pusat lingkaran sembilan titik dan titik Exeter. Pada segitiga
sama sisi keempat titik berimpit sehingga garis Euler tidak terdefinisi.

Parabola: tempat kedudukan titik-titik yang berjarak sama ke satu titik tetap (fokus) dan ke
satu garis tetap (direktriks). Pada contoh ini fokusnya C dan direktriksnya garis AB, sehingga
persamaannya: (jarak titik ke garis AB) = (jarak titik ke C). Karena direktriksnya miring,
parabolanya terputar (sumbunya tegak lurus AB).

\end{eulercomment}
\eulersubheading{Penyelesaian dengan EMT}
\begin{eulercomment}
Garis Euler adalah garis yang ditentukan dari sebarang segitiga yang tidak sama sisi. Garis ini
merupakan garis pusat segitiga, yang melalui beberapa titik penting yang ditentukan dari
segitiga itu, termasuk titik tinggi, pusat lingkaran luar, titik berat, titik Exeter, dan
pusat lingkaran sembilan titik segitiga tersebut.

Untuk demonstrasi, kita hitung dan gambar garis Euler pada suatu segitiga.

Pertama, kita definisikan titik-titik sudut segitiga di Euler. Kita memakai definisi yang
terlihat dalam ekspresi simbolik.
\end{eulercomment}
\begin{eulerprompt}
>A::=[-1,-1]; B::=[2,0]; C::=[1,2];
\end{eulerprompt}
\begin{eulercomment}
Untuk menggambar objek geometri, kita siapkan bidang gambar, dan tambahkan titik-titik ke
dalamnya. Semua gambar objek geometri ditambahkan ke gambar yang sedang aktif.
\end{eulercomment}
\begin{eulerprompt}
>setPlotRange(3); plotPoint(A,"A"); plotPoint(B,"B"); plotPoint(C,"C");
\end{eulerprompt}
\begin{eulercomment}
Kita juga dapat menambahkan sisi-sisi segitiga.
\end{eulercomment}
\begin{eulerprompt}
>plotSegment(A,B,""); plotSegment(B,C,""); plotSegment(C,A,""):
\end{eulerprompt}
\begin{eulercomment}
Berikut luas segitiga, dengan rumus determinan. Tentu saja, kita harus mengambil nilai
mutlak hasil ini.
\end{eulercomment}
\begin{eulerprompt}
>$areaTriangle(A,B,C)
\end{eulerprompt}
\begin{eulercomment}
Kita dapat menghitung koefisien sisi c.
\end{eulercomment}
\begin{eulerprompt}
>c &= lineThrough(A,B)
\end{eulerprompt}
\begin{euleroutput}
  
                              [- 1, 3, - 2]
  
\end{euleroutput}
\begin{eulercomment}
Dan juga mendapatkan rumus untuk garis ini.
\end{eulercomment}
\begin{eulerprompt}
>$getLineEquation(c,x,y)
\end{eulerprompt}
\begin{eulercomment}
Untuk bentuk Hesse, kita perlu menentukan suatu titik, sehingga titik tersebut berada pada sisi
positif bentuk Hesse. Memasukkan titik tersebut menghasilkan jarak positif ke garis.
\end{eulercomment}
\begin{eulerprompt}
>$getHesseForm(c,x,y,C), $at(%,[x=C[1],y=C[2]])
\end{eulerprompt}
\begin{eulercomment}
Sekarang kita hitung lingkaran luar ABC.
\end{eulercomment}
\begin{eulerprompt}
>LL &= circleThrough(A,B,C); $getCircleEquation(LL,x,y)
>O &= getCircleCenter(LL); $O
\end{eulerprompt}
\begin{eulercomment}
Gambar lingkaran dan pusatnya. Cu dan U bersifat simbolik. Kita evaluasi
ekspresi-ekspresi ini untuk Euler.
\end{eulercomment}
\begin{eulerprompt}
>plotCircle(LL()); plotPoint(O(),"O"):
\end{eulerprompt}
\begin{eulercomment}
Kita dapat menghitung perpotongan garis-garis tinggi di ABC (titik tinggi) secara numerik
dengan perintah berikut.
\end{eulercomment}
\begin{eulerprompt}
>H &= lineIntersection(perpendicular(A,lineThrough(C,B)),...
>  perpendicular(B,lineThrough(A,C))); $H
\end{eulerprompt}
\begin{eulercomment}
Sekarang kita dapat menghitung garis Euler segitiga tersebut.
\end{eulercomment}
\begin{eulerprompt}
>el &= lineThrough(H,O); $getLineEquation(el,x,y)
\end{eulerprompt}
\begin{eulercomment}
Tambahkan ke gambar kita.
\end{eulercomment}
\begin{eulerprompt}
>plotPoint(H(),"H"); plotLine(el(),"Garis Euler"):
\end{eulerprompt}
\begin{eulercomment}
Titik berat seharusnya berada pada garis ini.
\end{eulercomment}
\begin{eulerprompt}
>M &= (A+B+C)/3; $getLineEquation(el,x,y) with [x=M[1],y=M[2]]
>plotPoint(M(),"M"): // titik berat
\end{eulerprompt}
\begin{eulercomment}
Teori mengatakan MH=2*MO. Kita perlu menyederhanakan dengan radcan untuk mendapatkannya.
\end{eulercomment}
\begin{eulerprompt}
>$distance(M,H)/distance(M,O)|radcan
\end{eulerprompt}
\begin{eulercomment}
Terdapat juga fungsi untuk sudut.
\end{eulercomment}
\begin{eulerprompt}
>$computeAngle(A,C,B), degprint(%())
\end{eulerprompt}
\begin{euleroutput}
  60°15'18.43''
\end{euleroutput}
\begin{eulercomment}
Persamaan untuk pusat lingkaran dalam tidak terlalu bagus.
\end{eulercomment}
\begin{eulerprompt}
>Q &= lineIntersection(angleBisector(A,C,B),angleBisector(C,B,A))|radcan; $Q
\end{eulerprompt}
\begin{eulercomment}
Mari kita hitung juga ekspresi untuk jari-jari lingkaran dalam.
\end{eulercomment}
\begin{eulerprompt}
>r &= distance(Q,projectToLine(Q,lineThrough(A,B)))|ratsimp; $r
>LD &=  circleWithCenter(Q,r); // Lingkaran dalam
\end{eulerprompt}
\begin{eulercomment}
Mari kita tambahkan ini ke gambar.
\end{eulercomment}
\begin{eulerprompt}
>color(5); plotCircle(LD()):
\end{eulerprompt}
\begin{eulercomment}
Hasil yang diharapkan: luas = 7/2 (nilai mutlak); c: -x+3y=-2; O = [3/14,5/14];
H = [11/7,2/7]; garis Euler memuat M = [2/3,2/3] (kesamaan -1/2 = -1/2); MH/MO = 2;
sudut ACB = 60°15'18.43''.
Penjelasan: O, M, H segaris dan M membagi OH dengan perbandingan 1 : 2 (OM : MH = 1 : 2).

\end{eulercomment}
\eulersubheading{Parabola}
\begin{eulercomment}
Selanjutnya akan dicari persamaan tempat kedudukan titik-titik yang berjarak sama ke titik C
dan ke garis AB.
\end{eulercomment}
\begin{eulerprompt}
>p &= getHesseForm(lineThrough(A,B),x,y,C)-distance([x,y],C); $p='0
\end{eulerprompt}
\begin{eulercomment}
Persamaan tersebut dapat digambar menjadi satu dengan gambar sebelumnya.
\end{eulercomment}
\begin{eulerprompt}
>plot2d(p,level=0,add=1,contourcolor=6):
\end{eulerprompt}
\begin{eulercomment}
Ini seharusnya berupa suatu fungsi, tetapi penyelesai (solver) standar Maxima hanya dapat
menemukan solusinya jika persamaan dikuadratkan. Akibatnya, kita mendapatkan solusi palsu.
\end{eulercomment}
\begin{eulerprompt}
>akar &= solve(getHesseForm(lineThrough(A,B),x,y,C)^2-distance([x,y],C)^2,y)
\end{eulerprompt}
\begin{euleroutput}
  
          [y = - 3 x - sqrt(70) sqrt(9 - 2 x) + 26, 
                                y = - 3 x + sqrt(70) sqrt(9 - 2 x) + 26]
  
\end{euleroutput}
\begin{eulercomment}
Solusi pertama adalah

maxima: akar[1]

Menambahkan solusi pertama ke gambar menunjukkan bahwa itu memang lintasan yang kita cari.
Teori mengatakan bahwa itu adalah parabola yang terputar.
\end{eulercomment}
\begin{eulerprompt}
>plot2d(&rhs(akar[1]),add=1):
>function g(x) &= rhs(akar[1]); $'g(x)= g(x)// fungsi yang mendefinisikan kurva di atas
>T &=[-1, g(-1)]; // ambil sebarang titik pada kurva tersebut
>dTC &= distance(T,C); $fullratsimp(dTC), $float(%) // jarak T ke C
>U &= projectToLine(T,lineThrough(A,B)); $U // proyeksi T pada garis AB
>dU2AB &= distance(T,U); $fullratsimp(dU2AB), $float(%) // jarak T ke AB
\end{eulerprompt}
\begin{eulercomment}
Ternyata jarak T ke C sama dengan jarak T ke AB. Coba Anda pilih titik T yang lain dan
ulangi perhitungan-perhitungan di atas untuk menunjukkan bahwa hasilnya juga sama.

*** Latihan Contoh 4: Memilih titik T yang lain pada parabola

Titik T diambil pada x = 0, x = 2, dan x = 4 (daerah asal kurva: 9 - 2x \textgreater{}= 0, jadi x \textless{}= 4.5).
Jarak ke fokus C dan jarak ke direktriks AB dihitung secara numerik.
\end{eulercomment}
\begin{eulerprompt}
>T2=[0,g(0)]; distance(T2,C), jarakKeGaris(T2,lineThrough(A,B))
\end{eulerprompt}
\begin{euleroutput}
  1.48645948177
  1.48645948177
\end{euleroutput}
\begin{eulerprompt}
>T3=[2,g(2)]; distance(T3,C), jarakKeGaris(T3,lineThrough(A,B))
\end{eulerprompt}
\begin{euleroutput}
  1.22542661171
  1.22542661171
\end{euleroutput}
\begin{eulerprompt}
>T4=[4,g(4)]; distance(T4,C), jarakKeGaris(T4,lineThrough(A,B))
\end{eulerprompt}
\begin{euleroutput}
  4.71185670748
  4.71185670748
\end{euleroutput}
\begin{eulercomment}
Hasil yang diharapkan (jarak ke C, jarak ke AB):
x = 0: (1.48646, 1.48646); x = 2: (1.22543, 1.22543); x = 4: (4.71186, 4.71186).
Penjelasan: untuk setiap titik T pada kurva, jaraknya ke C sama dengan jaraknya ke garis AB.
Inilah sifat definisi parabola (jarak ke fokus = jarak ke direktriks). Titik di x=-1 memberi
2.1356 untuk keduanya.

\begin{eulercomment}
\eulerheading{Contoh 5: Trigonometri Rasional}
\begin{eulercomment}
\end{eulercomment}
\eulersubheading{Uraian Materi}
\begin{eulercomment}
Trigonometri rasional (N. J. Wildberger, "Divine Proportions") mengganti jarak dengan kuadran
(quadrance) dan sudut dengan spread, sehingga banyak perhitungan cukup memakai bilangan
rasional.

- Kuadran dua titik: Q(A,B) = (x2-x1)\textasciicircum{}2 + (y2-y1)\textasciicircum{}2 (kuadrat jarak).
- Spread dua garis: s = sin\textasciicircum{}2(theta), nilainya 0 (sejajar) sampai 1 (tegak lurus).
- Teorema Pythagoras menjadi Q1 + Q2 = Q3 (segitiga siku-siku, spread 1).
- Hukum silang (cross law): (Q1+Q2-Q3)\textasciicircum{}2 = 4 Q1 Q2 (1 - s3), setara hukum kosinus.
- Hukum spread: s1/Q1 = s2/Q2 = s3/Q3, setara hukum sinus.
- Spread rangkap tiga (triple spread): (s1+s2+s3)\textasciicircum{}2 = 2(s1\textasciicircum{}2+s2\textasciicircum{}2+s3\textasciicircum{}2) + 4 s1 s2 s3.

Perhitungan di bawah memakai Maxima, sehingga keuntungan utama trigonometri rasional
(dapat dikerjakan dengan kertas dan pensil saja) tidak terlihat; Anda dianjurkan memeriksa
hasil-hasilnya tanpa komputer.

\end{eulercomment}
\eulersubheading{Penyelesaian dengan EMT}
\begin{eulercomment}
Ini terinspirasi dari ceramah N.J.Wildberger. Dalam bukunya "Divine Proportions", Wildberger
mengusulkan penggantian pengertian klasik jarak dan sudut dengan kuadran dan spread.
Dengan keduanya, memang dimungkinkan menghindari fungsi trigonometri pada banyak contoh, dan
tetap "rasional".

Berikut ini, saya memperkenalkan konsep-konsepnya, dan menyelesaikan beberapa soal. Saya memakai
perhitungan simbolik Maxima di sini, yang menyembunyikan keuntungan utama trigonometri rasional,
yaitu bahwa perhitungan dapat dilakukan hanya dengan kertas dan pensil. Anda diundang untuk
memeriksa hasil-hasilnya tanpa komputer.

Intinya, perhitungan rasional simbolik sering menghasilkan hasil yang sederhana. Sebaliknya,
trigonometri klasik menghasilkan hasil trigonometri yang rumit, yang hanya bisa dievaluasi
menjadi hampiran numerik.

Catatan perbaikan: variabel Maxima dibersihkan dulu karena cc, a, b, c dari contoh sebelumnya
masih tersimpan dan membuat hasil crosslaw di versi asli salah.
\end{eulercomment}
\begin{eulerprompt}
>&remvalue();
>load geometry;
\end{eulerprompt}
\begin{eulercomment}
Sebagai pengenalan pertama, kita pakai segitiga siku-siku dengan perbandingan Mesir yang
terkenal 3, 4 dan 5. Perintah-perintah berikut adalah perintah Euler untuk menggambar
geometri bidang yang terdapat dalam file Euler "geometry.e".
\end{eulercomment}
\begin{eulerprompt}
>C&:=[0,0]; A&:=[4,0]; B&:=[0,3]; ...
>setPlotRange(-1,5,-1,5); ...
>plotPoint(A,"A"); plotPoint(B,"B"); plotPoint(C,"C"); ...
>plotSegment(B,A,"c"); plotSegment(A,C,"b"); plotSegment(C,B,"a"); ...
>insimg(30);
\end{eulerprompt}
\begin{eulercomment}
Tentu saja,

\end{eulercomment}
\begin{eulerformula}
\[
\sin(w_a)=\frac{a}{c},
\]
\end{eulerformula}
\begin{eulercomment}
dengan wa adalah sudut di A. Cara biasa untuk menghitung sudut ini adalah dengan mengambil
invers fungsi sinus. Hasilnya adalah sudut yang "tidak enak dicerna", yang hanya bisa
dicetak secara hampiran.
\end{eulercomment}
\begin{eulerprompt}
>wa := arcsin(3/5); degprint(wa)
\end{eulerprompt}
\begin{euleroutput}
  36°52'11.63''
\end{euleroutput}
\begin{eulercomment}
Trigonometri rasional mencoba menghindari hal ini.

Pengertian pertama trigonometri rasional adalah kuadran, yang menggantikan jarak. Sebenarnya,
kuadran hanyalah jarak yang dikuadratkan. Berikut ini, a, b, dan c menyatakan kuadran
sisi-sisinya.

Teorema Pythogoras kemudian menjadi sederhana: a+b=c.
\end{eulercomment}
\begin{eulerprompt}
>a &= 3^2; b &= 4^2; c &= 5^2; &a+b=c
\end{eulerprompt}
\begin{euleroutput}
  
                                 25 = 25
  
\end{euleroutput}
\begin{eulercomment}
Pengertian kedua trigonometri rasional adalah spread. Spread mengukur bukaan antara
garis-garis. Nilainya 0, jika garis-garisnya sejajar, dan 1, jika garis-garisnya tegak lurus.
Spread adalah kuadrat sinus sudut antara dua garis.

Spread garis AB dan AC pada gambar di atas didefinisikan sebagai

\end{eulercomment}
\begin{eulerformula}
\[
s_a = \sin(\alpha)^2 = \frac{a}{c},
\]
\end{eulerformula}
\begin{eulercomment}
dengan a dan c adalah kuadran-kuadran pada sebarang segitiga siku-siku dengan satu sudut di A.
\end{eulercomment}
\begin{eulerprompt}
>sa &= a/c; $sa
\end{eulerprompt}
\begin{eulercomment}
Ini tentunya lebih mudah dihitung daripada sudut. Tetapi Anda kehilangan sifat bahwa
sudut-sudut dapat dijumlahkan dengan mudah.

Tentu saja, kita dapat mengubah nilai hampiran sudut wa menjadi spread, dan mencetaknya
sebagai pecahan.
\end{eulercomment}
\begin{eulerprompt}
>fracprint(sin(wa)^2)
\end{eulerprompt}
\begin{euleroutput}
  9/25
\end{euleroutput}
\begin{eulercomment}
Hukum kosinus pada trigonometri klasik diterjemahkan menjadi "hukum silang" berikut.

\end{eulercomment}
\begin{eulerformula}
\[
(c+b-a)^2 = 4 b c \, (1-s_a)
\]
\end{eulerformula}
\begin{eulercomment}
Di sini a, b, dan c adalah kuadran sisi-sisi segitiga, dan sa adalah spread di titik sudut A.
Sisi a, seperti biasa, berhadapan dengan titik sudut A.

Hukum-hukum ini sudah diimplementasikan dalam file geometry.e yang kita panggil ke Euler.
\end{eulercomment}
\begin{eulerprompt}
>$crosslaw(aa,bb,cc,saa)
\end{eulerprompt}
\begin{eulercomment}
Pada kasus kita didapat
\end{eulercomment}
\begin{eulerprompt}
>$crosslaw(a,b,c,sa)
\end{eulerprompt}
\begin{eulercomment}
Mari kita pakai hukum silang ini untuk mencari spread di A. Untuk itu, kita buat hukum silang
untuk kuadran a, b, dan c, lalu selesaikan untuk spread sa yang tidak diketahui.

Anda dapat melakukan ini dengan mudah secara manual, tetapi saya memakai Maxima. Tentu saja,
kita mendapatkan hasil yang sudah kita punya.
\end{eulercomment}
\begin{eulerprompt}
>$crosslaw(a,b,c,x), $solve(%,x)
\end{eulerprompt}
\begin{eulercomment}
Kita sudah mengetahui ini. Definisi spread adalah kasus khusus dari hukum silang.

Kita juga dapat menyelesaikan ini untuk a,b,c umum. Hasilnya adalah rumus yang menghitung
spread suatu sudut segitiga bila diketahui kuadran ketiga sisinya.
\end{eulercomment}
\begin{eulerprompt}
>$solve(crosslaw(aa,bb,cc,x),x)
\end{eulerprompt}
\begin{eulercomment}
Kita dapat membuat fungsi dari hasil itu. Fungsi seperti itu sudah didefinisikan dalam file
geometry.e milik Euler.
\end{eulercomment}
\begin{eulerprompt}
>$spread(a,b,c)
\end{eulerprompt}
\begin{eulercomment}
Sebagai contoh, kita dapat memakainya untuk menghitung sudut segitiga dengan sisi-sisi

\end{eulercomment}
\begin{eulerformula}
\[
a, \quad a, \quad \frac{4a}{7}
\]
\end{eulerformula}
\begin{eulercomment}
Hasilnya rasional, yang tidak mudah didapat bila kita memakai trigonometri klasik.
\end{eulercomment}
\begin{eulerprompt}
>$spread(a,a,4*a/7)
\end{eulerprompt}
\begin{eulercomment}
Ini adalah sudut dalam derajat.
\end{eulercomment}
\begin{eulerprompt}
>degprint(arcsin(sqrt(6/7)))
\end{eulerprompt}
\begin{euleroutput}
  67°47'32.44''
\end{euleroutput}
\begin{eulercomment}
\end{eulercomment}
\eulersubheading{Contoh Lain}
\begin{eulercomment}
Sekarang, mari kita coba contoh yang lebih lanjut.

Kita tentukan tiga titik sudut segitiga sebagai berikut.
\end{eulercomment}
\begin{eulerprompt}
>A&:=[1,2]; B&:=[4,3]; C&:=[0,4]; ...
>setPlotRange(-1,5,1,7); ...
>plotPoint(A,"A"); plotPoint(B,"B"); plotPoint(C,"C"); ...
>plotSegment(B,A,"c"); plotSegment(A,C,"b"); plotSegment(C,B,"a"); ...
>insimg;
\end{eulerprompt}
\begin{eulercomment}
Dengan Pythogoras, mudah menghitung jarak antara dua titik. Mula-mula saya pakai fungsi
distance dari file Euler untuk geometri. Fungsi distance memakai geometri klasik.
\end{eulercomment}
\begin{eulerprompt}
>$distance(A,B)
\end{eulerprompt}
\begin{eulercomment}
Euler juga memiliki fungsi untuk kuadran antara dua titik.

Pada contoh berikut, karena c+b bukan a, segitiga ini bukan segitiga siku-siku.
\end{eulercomment}
\begin{eulerprompt}
>c &= quad(A,B); $c, b &= quad(A,C); $b, a &= quad(B,C); $a,
\end{eulerprompt}
\begin{eulercomment}
Pertama, mari kita hitung sudut tradisional. Fungsi computeAngle memakai metode biasa yang
berdasarkan hasil kali titik dua vektor. Hasilnya berupa hampiran bilangan floating point.

\end{eulercomment}
\begin{eulerformula}
\[
A=<1,2>\quad B=<4,3>,\quad C=<0,4>
\]
\end{eulerformula}
\begin{eulerformula}
\[
\mathbf{a}=C-B=<-4,1>,\quad \mathbf{c}=A-B=<-3,-1>,\quad \beta=\angle ABC
\]
\end{eulerformula}
\begin{eulerformula}
\[
\mathbf{a}.\mathbf{c}=|\mathbf{a}|.|\mathbf{c}|\cos \beta
\]
\end{eulerformula}
\begin{eulerformula}
\[
\cos \angle ABC =\cos\beta=\frac{\mathbf{a}.\mathbf{c}}{|\mathbf{a}|.|\mathbf{c}|}=\frac{12-1}{\sqrt{17}\sqrt{10}}=\frac{11}{\sqrt{17}\sqrt{10}}
\]
\end{eulerformula}
\begin{eulerprompt}
>wb &= computeAngle(A,B,C); $wb, $(wb/pi*180)()
\end{eulerprompt}
\begin{euleroutput}
  32.4711922908
\end{euleroutput}
\begin{eulercomment}
Dengan pensil dan kertas, kita dapat melakukan hal yang sama dengan hukum silang. Kita
masukkan kuadran a, b, dan c ke dalam hukum silang dan selesaikan untuk x.
\end{eulercomment}
\begin{eulerprompt}
>$crosslaw(a,b,c,x), $solve(%,x), //(b+c-a)^=4b.c(1-x)
\end{eulerprompt}
\begin{eulercomment}
Itulah yang dikerjakan fungsi spread yang didefinisikan dalam "geometry.e".
\end{eulercomment}
\begin{eulerprompt}
>sb &= spread(b,a,c); $sb
\end{eulerprompt}
\begin{eulercomment}
Maxima mendapatkan hasil yang sama dengan trigonometri biasa, bila kita paksa. Ia
menyelesaikan suku sin(arccos(...)) menjadi hasil pecahan. Kebanyakan mahasiswa tidak
dapat melakukan ini.
\end{eulercomment}
\begin{eulerprompt}
>$sin(computeAngle(A,B,C))^2
\end{eulerprompt}
\begin{eulercomment}
Setelah kita punya spread di B, kita dapat menghitung tinggi ha pada sisi a. Ingat bahwa

\end{eulercomment}
\begin{eulerformula}
\[
s_b=\frac{h_a}{c}
\]
\end{eulerformula}
\begin{eulercomment}
menurut definisi.
\end{eulercomment}
\begin{eulerprompt}
>ha &= c*sb; $ha
\end{eulerprompt}
\begin{eulercomment}
Gambar berikut dihasilkan dengan program geometri C.a.R., yang dapat menggambar kuadran dan
spread.

Menurut definisi, panjang ha adalah akar kuadrat dari kuadrannya.
\end{eulercomment}
\begin{eulerprompt}
>$sqrt(ha)
\end{eulerprompt}
\begin{eulercomment}
Sekarang kita dapat menghitung luas segitiga. Jangan lupa, bahwa kita sedang berurusan
dengan kuadran!
\end{eulercomment}
\begin{eulerprompt}
>$sqrt(ha)*sqrt(a)/2
\end{eulerprompt}
\begin{eulercomment}
Rumus determinan yang biasa menghasilkan hal yang sama.
\end{eulercomment}
\begin{eulerprompt}
>$areaTriangle(B,A,C)
\end{eulerprompt}
\begin{eulercomment}
Hasil yang diharapkan: c = 10, b = 5, a = 17; sudut B = 32.4712°; sb = 49/170; ha = 49/17;
luas = 7/2 (sama dengan rumus determinan).

\end{eulercomment}
\eulersubheading{Rumus Heron}
\begin{eulercomment}
Sekarang, mari kita selesaikan soal ini secara umum!
\end{eulercomment}
\begin{eulerprompt}
>&remvalue(a,b,c,sb,ha);
\end{eulerprompt}
\begin{eulercomment}
Pertama kita hitung spread di B untuk segitiga dengan sisi a, b, dan c. Kemudian kita hitung
luas kuadrat (disebut "quadrea"?), memfaktorkannya dengan Maxima, dan kita mendapatkan rumus
Heron yang terkenal.

Harus diakui, ini sulit dikerjakan dengan pensil dan kertas.
\end{eulercomment}
\begin{eulerprompt}
>$spread(b^2,c^2,a^2), $factor(%*c^2*a^2/4)
\end{eulerprompt}
\begin{eulercomment}
\end{eulercomment}
\eulersubheading{Aturan Spread Rangkap Tiga}
\begin{eulercomment}
Kelemahan spread adalah bahwa spread tidak lagi dijumlahkan seperti sudut.

Namun, ketiga spread suatu segitiga memenuhi aturan "spread rangkap tiga" berikut.
\end{eulercomment}
\begin{eulerprompt}
>&remvalue(sa,sb,sc); $triplespread(sa,sb,sc)
\end{eulerprompt}
\begin{eulercomment}
Aturan ini berlaku untuk sebarang tiga sudut yang berjumlah 180°.

\end{eulercomment}
\begin{eulerformula}
\[
\alpha+\beta+\gamma=\pi
\]
\end{eulerformula}
\begin{eulercomment}
Karena spread

\end{eulercomment}
\begin{eulerformula}
\[
\alpha, \pi-\alpha
\]
\end{eulerformula}
\begin{eulercomment}
sama, aturan spread rangkap tiga juga benar, jika

\end{eulercomment}
\begin{eulerformula}
\[
\alpha+\beta=\gamma
\]
\end{eulerformula}
\begin{eulercomment}
Karena spread sudut negatif sama, aturan spread rangkap tiga juga berlaku, jika

\end{eulercomment}
\begin{eulerformula}
\[
\alpha+\beta+\gamma=0
\]
\end{eulerformula}
\begin{eulercomment}
Sebagai contoh, kita dapat menghitung spread sudut 60°. Nilainya 3/4. Persamaan ini punya
solusi kedua, yaitu semua spread bernilai 0.
\end{eulercomment}
\begin{eulerprompt}
>$solve(triplespread(x,x,x),x)
\end{eulerprompt}
\begin{eulercomment}
Spread sudut 90° jelas 1. Jika dua sudut berjumlah 90°, spread keduanya menyelesaikan
persamaan spread rangkap tiga dengan a,b,1. Dengan perhitungan berikut kita mendapatkan a+b=1.
\end{eulercomment}
\begin{eulerprompt}
>$triplespread(x,y,1), $solve(%,x)
\end{eulerprompt}
\begin{eulercomment}
Karena spread 180°-t sama dengan spread t, rumus spread rangkap tiga juga berlaku, jika satu
sudut adalah jumlah atau selisih dua sudut lainnya.

Jadi kita dapat mencari spread sudut rangkap. Perhatikan bahwa ada dua solusi lagi. Kita
jadikan ini sebuah fungsi.
\end{eulercomment}
\begin{eulerprompt}
>$solve(triplespread(a,a,x),x), function doublespread(a) &= factor(rhs(%[1]))
\end{eulerprompt}
\begin{euleroutput}
  
                              - 4 (a - 1) a
  
\end{euleroutput}
\begin{eulercomment}
Hasil yang diharapkan: spread 60° = 3/4; a+b = 1 untuk dua sudut komplementer;
doublespread(a) = 4a(1-a).

\end{eulercomment}
\eulersubheading{Garis Bagi Sudut}
\begin{eulercomment}
Ini adalah situasi yang sudah kita ketahui.
\end{eulercomment}
\begin{eulerprompt}
>C&:=[0,0]; A&:=[4,0]; B&:=[0,3]; ...
>setPlotRange(-1,5,-1,5); ...
>plotPoint(A,"A"); plotPoint(B,"B"); plotPoint(C,"C"); ...
>plotSegment(B,A,"c"); plotSegment(A,C,"b"); plotSegment(C,B,"a"); ...
>insimg;
\end{eulerprompt}
\begin{eulercomment}
Mari kita hitung panjang garis bagi sudut di A. Tetapi kita ingin menyelesaikannya untuk
a,b,c umum.
\end{eulercomment}
\begin{eulerprompt}
>&remvalue(a,b,c);
\end{eulerprompt}
\begin{eulercomment}
Jadi pertama-tama kita hitung spread sudut yang dibagi dua di A, memakai rumus spread
rangkap tiga.

Masalah rumus ini muncul lagi. Rumus ini punya dua solusi. Kita harus memilih yang benar.
Solusi yang lain berkaitan dengan sudut yang dibagi dua 180°-wa.
\end{eulercomment}
\begin{eulerprompt}
>$triplespread(x,x,a/(a+b)), $solve(%,x), sa2 &= rhs(%[1]); $sa2
\end{eulerprompt}
\begin{eulercomment}
Mari kita periksa untuk segiempat (segitiga) Mesir.
\end{eulercomment}
\begin{eulerprompt}
>$sa2 with [a=3^2,b=4^2]
\end{eulerprompt}
\begin{eulercomment}
Kita dapat mencetak sudutnya di Euler, setelah memindahkan spread ke radian.
\end{eulercomment}
\begin{eulerprompt}
>wa2 := arcsin(sqrt(1/10)); degprint(wa2)
\end{eulerprompt}
\begin{euleroutput}
  18°26'5.82''
\end{euleroutput}
\begin{eulercomment}
Titik P adalah perpotongan garis bagi sudut dengan sumbu-y.
\end{eulercomment}
\begin{eulerprompt}
>P := [0,tan(wa2)*4]
\end{eulerprompt}
\begin{euleroutput}
  [0,  1.33333]
\end{euleroutput}
\begin{eulerprompt}
>plotPoint(P,"P"); plotSegment(A,P):
\end{eulerprompt}
\begin{eulercomment}
Mari kita periksa sudut-sudut pada contoh khusus kita.
\end{eulercomment}
\begin{eulerprompt}
>computeAngle(C,A,P), computeAngle(P,A,B)
\end{eulerprompt}
\begin{euleroutput}
  0.321750554397
  0.321750554397
\end{euleroutput}
\begin{eulercomment}
Sekarang kita hitung panjang garis bagi AP.

Kita pakai teorema sinus pada segitiga APC. Teorema ini menyatakan bahwa

\end{eulercomment}
\begin{eulerformula}
\[
\frac{BC}{\sin(w_a)} = \frac{AC}{\sin(w_b)} = \frac{AB}{\sin(w_c)}
\]
\end{eulerformula}
\begin{eulercomment}
berlaku pada sebarang segitiga. Kuadratkan, maka ia diterjemahkan menjadi apa yang disebut
"hukum spread"

\end{eulercomment}
\begin{eulerformula}
\[
\frac{a}{s_a} = \frac{b}{s_b} = \frac{c}{s_b}
\]
\end{eulerformula}
\begin{eulercomment}
dengan a,b,c menyatakan kuadran-kuadran.

Karena spread CPA adalah 1-sa2, dari situ kita dapatkan bisa/1=b/(1-sa2) dan dapat menghitung
bisa (kuadran garis bagi sudut).
\end{eulercomment}
\begin{eulerprompt}
>&factor(ratsimp(b/(1-sa2))); bisa &= %; $bisa
\end{eulerprompt}
\begin{eulercomment}
Mari kita periksa rumus ini untuk nilai-nilai Mesir kita.
\end{eulercomment}
\begin{eulerprompt}
>sqrt(mxmeval("at(bisa,[a=3^2,b=4^2])")), distance(A,P)
\end{eulerprompt}
\begin{euleroutput}
  4.21637021356
  4.21637021356
\end{euleroutput}
\begin{eulercomment}
Kita juga dapat menghitung P memakai rumus spread.
\end{eulercomment}
\begin{eulerprompt}
>py&=factor(ratsimp(sa2*bisa)); $py
\end{eulerprompt}
\begin{eulercomment}
Nilainya sama dengan yang kita dapat dengan rumus trigonometri.
\end{eulercomment}
\begin{eulerprompt}
>sqrt(mxmeval("at(py,[a=3^2,b=4^2])"))
\end{eulerprompt}
\begin{euleroutput}
  1.33333333333
\end{euleroutput}
\begin{eulercomment}
Hasil yang diharapkan: sa2 = 1/10; sudut 18°26'5.82''; P = [0, 1.33333];
AP = 4.21637; kedua sudut CAP dan PAB = 0.321751 rad.

\end{eulercomment}
\eulersubheading{Sudut Tali Busur (Chord Angle)}
\begin{eulercomment}
Perhatikan situasi berikut.
\end{eulercomment}
\begin{eulerprompt}
>setPlotRange(1.2); ...
>color(1); plotCircle(circleWithCenter([0,0],1)); ...
>A:=[cos(1),sin(1)]; B:=[cos(2),sin(2)]; C:=[cos(6),sin(6)]; ...
>plotPoint(A,"A"); plotPoint(B,"B"); plotPoint(C,"C"); ...
>color(3); plotSegment(A,B,"c"); plotSegment(A,C,"b"); plotSegment(C,B,"a"); ...
>color(1); O:=[0,0];  plotPoint(O,"0"); ...
>plotSegment(A,O); plotSegment(B,O); plotSegment(C,O,"r"); ...
>insimg;
\end{eulerprompt}
\begin{eulercomment}
Kita dapat memakai Maxima untuk menyelesaikan rumus spread rangkap tiga untuk sudut-sudut di
pusat O terhadap r. Dengan demikian kita mendapatkan rumus untuk kuadrat jari-jari
lingkaran luar dalam bentuk kuadran-kuadran sisinya.

Kali ini, Maxima menghasilkan beberapa nol kompleks, yang kita abaikan.
\end{eulercomment}
\begin{eulerprompt}
>&remvalue(a,b,c,r); // hapus nilai-nilai sebelumnya untuk perhitungan baru
>rabc &= rhs(solve(triplespread(spread(b,r,r),spread(a,r,r),spread(c,r,r)),r)[4]); $rabc
\end{eulerprompt}
\begin{eulercomment}
Kita dapat menjadikan itu fungsi Euler.
\end{eulercomment}
\begin{eulerprompt}
>function periradius(a,b,c) &= rabc;
\end{eulerprompt}
\begin{eulercomment}
Mari kita periksa hasilnya untuk titik-titik A,B,C kita.
\end{eulercomment}
\begin{eulerprompt}
>a:=quadrance(B,C); b:=quadrance(A,C); c:=quadrance(A,B);
\end{eulerprompt}
\begin{eulercomment}
Jari-jarinya memang 1.
\end{eulercomment}
\begin{eulerprompt}
>periradius(a,b,c)
\end{eulerprompt}
\begin{euleroutput}
  1
\end{euleroutput}
\begin{eulercomment}
Faktanya, spread CBA hanya bergantung pada b dan c. Ini adalah teorema sudut tali busur.
\end{eulercomment}
\begin{eulerprompt}
>$spread(b,a,c)*rabc | ratsimp
\end{eulerprompt}
\begin{eulercomment}
Sebenarnya spread tersebut adalah b/(4r), dan kita lihat bahwa sudut tali busur dari tali
busur b adalah setengah sudut pusat.
\end{eulercomment}
\begin{eulerprompt}
>$doublespread(b/(4*r))-spread(b,r,r) | ratsimp
\end{eulerprompt}
\begin{eulercomment}
Hasil yang diharapkan: periradius = 1; hasil kali = b/4; selisih terakhir = 0.

\begin{eulercomment}
\eulerheading{Contoh 6: Jarak Minimal pada Bidang}
\begin{eulercomment}
\end{eulercomment}
\eulersubheading{Uraian Materi}
\begin{eulercomment}
Titik Fermat (titik Torricelli) suatu segitiga adalah titik F yang meminimumkan jumlah jarak
FA + FB + FC. Jika semua sudut segitiga kurang dari 120 derajat, F berada di dalam segitiga
dan melihat setiap sisi dengan sudut 120 derajat (sudut AFB = BFC = CFA = 120°). Jika ada
sudut yang sama atau lebih dari 120 derajat, titik minimumnya adalah titik sudut tersebut.
Konstruksinya: buat segitiga sama sisi di luar pada dua sisi; garis dari titik sudut ke puncak
segitiga sama sisi di hadapannya berpotongan di F.

Untuk dua titik A dan B, himpunan titik dengan MA+MB = konstanta adalah elips dengan fokus A dan B.

\end{eulercomment}
\eulersubheading{Penyelesaian dengan EMT}
\begin{eulercomment}
\end{eulercomment}
\eulersubheading{Catatan pendahuluan}
\begin{eulercomment}
Fungsi yang mengaitkan titik M pada bidang dengan jarak AM antara titik tetap A dan M,
mempunyai garis-garis tinggi (kontur) yang cukup sederhana: lingkaran-lingkaran dengan pusat A.
\end{eulercomment}
\begin{eulerprompt}
>&remvalue();
>A=[-1,-1];
>function d1(x,y):=sqrt((x-A[1])^2+(y-A[2])^2)
>fcontour("d1",xmin=-2,xmax=0,ymin=-2,ymax=0,hue=1, ...
>title="Jika tampak elips, atur jendela Anda menjadi persegi"):
\end{eulerprompt}
\begin{eulercomment}
dan grafiknya juga cukup sederhana: bagian atas sebuah kerucut:
\end{eulercomment}
\begin{eulerprompt}
>plot3d("d1",xmin=-2,xmax=0,ymin=-2,ymax=0):
\end{eulerprompt}
\begin{eulercomment}
Tentu saja minimum 0 dicapai di A.

\end{eulercomment}
\eulersubheading{Dua titik}
\begin{eulercomment}
Sekarang kita lihat fungsi MA+MB dengan A dan B dua titik (tetap). Ada "fakta yang terkenal"
bahwa kurva-kurva kontur-nya berupa elips, dengan titik fokus A dan B; kecuali minimum AB yang
konstan pada ruas [AB]:
\end{eulercomment}
\begin{eulerprompt}
>B=[1,-1];
>function d2(x,y):=d1(x,y)+sqrt((x-B[1])^2+(y-B[2])^2)
>fcontour("d2",xmin=-2,xmax=2,ymin=-3,ymax=1,hue=1):
\end{eulerprompt}
\begin{eulercomment}
Grafiknya lebih menarik:
\end{eulercomment}
\begin{eulerprompt}
>plot3d("d2",xmin=-2,xmax=2,ymin=-3,ymax=1):
\end{eulerprompt}
\begin{eulercomment}
Pembatasan pada garis (AB) lebih terkenal:
\end{eulercomment}
\begin{eulerprompt}
>plot2d("abs(x+1)+abs(x-1)",xmin=-3,xmax=3):
\end{eulerprompt}
\begin{eulercomment}
\end{eulercomment}
\eulersubheading{Tiga titik}
\begin{eulercomment}
Sekarang hal-hal menjadi kurang sederhana: Agak kurang terkenal bahwa MA+MB+MC mencapai
minimumnya di satu titik pada bidang, tetapi menentukannya tidak sesederhana itu:

1) Jika salah satu sudut segitiga ABC lebih dari 120° (misalkan di A), maka minimum dicapai
tepat di titik ini (misalkan AB+AC).

Contoh:
\end{eulercomment}
\begin{eulerprompt}
>C=[-4,1];
>function d3(x,y):=d2(x,y)+sqrt((x-C[1])^2+(y-C[2])^2)
>plot3d("d3",xmin=-5,xmax=3,ymin=-4,ymax=4);
\end{eulerprompt}
\begin{euleroutput}
  Illegal function result in map.
   %ploteval2:
      return map(f$,x,y;args());
  f3dplot:
      z=%ploteval2(f$,x,y,maps;args());
  Try "trace errors" to inspect local variables after errors.
  plot3d:
      =cp,=cplevel,=cpcolor,=cpdelta);
\end{euleroutput}
\begin{eulerprompt}
>insimg;
>fcontour("d3",xmin=-4,xmax=1,ymin=-2,ymax=2,hue=1,title="Minimum ada di A");
\end{eulerprompt}
\begin{euleroutput}
  Illegal function result in map.
   %ploteval2:
      return map(f$,x,y;args());
  Try "trace errors" to inspect local variables after errors.
  fcontour:
      Z=%ploteval2(f$,X,Y,maps;args());
\end{euleroutput}
\begin{eulerprompt}
>P=(A_B_C_A)'; plot2d(P[1],P[2],add=1,color=12);
>insimg;
\end{eulerprompt}
\begin{eulercomment}
2) Tetapi jika semua sudut segitiga ABC kurang dari 120°, minimum berada pada suatu titik F di
bagian dalam segitiga, yang merupakan satu-satunya titik yang melihat sisi-sisi ABC dengan
sudut yang sama (maka masing-masing 120°):
\end{eulercomment}
\begin{eulerprompt}
>C=[-0.5,1];
>plot3d("d3",xmin=-2,xmax=2,ymin=-2,ymax=2):
\end{eulerprompt}
\begin{euleroutput}
  Illegal function result in map.
   %ploteval2:
      return map(f$,x,y;args());
  f3dplot:
      z=%ploteval2(f$,x,y,maps;args());
  Try "trace errors" to inspect local variables after errors.
  plot3d:
      =cp,=cplevel,=cpcolor,=cpdelta);
\end{euleroutput}
\begin{eulerprompt}
>fcontour("d3",xmin=-2,xmax=2,ymin=-2,ymax=2,hue=1,title="Titik Fermat");
\end{eulerprompt}
\begin{euleroutput}
  Illegal function result in map.
   %ploteval2:
      return map(f$,x,y;args());
  Try "trace errors" to inspect local variables after errors.
  fcontour:
      Z=%ploteval2(f$,X,Y,maps;args());
\end{euleroutput}
\begin{eulerprompt}
>P=(A_B_C_A)'; plot2d(P[1],P[2],add=1,color=12);
>insimg;
\end{eulerprompt}
\begin{eulercomment}
Merupakan kegiatan yang menarik untuk mewujudkan gambar di atas dengan perangkat lunak
geometri; misalnya, saya tahu satu perangkat lunak yang ditulis dalam Java dan memiliki
instruksi "garis kontur"...

Semua hal di atas ditemukan oleh seorang hakim Perancis bernama Pierre de Fermat; ia menulis
surat kepada para dilettante lain seperti pastor Marin Mersenne dan Blaise Pascal yang bekerja
di bagian pajak penghasilan. Jadi titik tunggal F sehingga FA+FB+FC minimal disebut titik
Fermat segitiga tersebut. Tetapi tampaknya beberapa tahun sebelumnya, orang Italia Torricelli
telah menemukan titik ini sebelum Fermat! Bagaimanapun tradisinya adalah menandai titik ini F...

\end{eulercomment}
\eulersubheading{Empat titik}
\begin{eulercomment}
Langkah berikutnya adalah menambahkan titik keempat D dan mencoba meminimumkan MA+MB+MC+MD;
misalkan Anda adalah operator TV kabel dan ingin mencari di lahan mana antena harus dipasang
agar Anda bisa melayani empat desa dan memakai kabel sependek mungkin!
\end{eulercomment}
\begin{eulerprompt}
>D=[1,1];
>function d4(x,y):=d3(x,y)+sqrt((x-D[1])^2+(y-D[2])^2)
>plot3d("d4",xmin=-1.5,xmax=1.5,ymin=-1.5,ymax=1.5):
\end{eulerprompt}
\begin{euleroutput}
  Illegal function result in map.
   %ploteval2:
      return map(f$,x,y;args());
  f3dplot:
      z=%ploteval2(f$,x,y,maps;args());
  Try "trace errors" to inspect local variables after errors.
  plot3d:
      =cp,=cplevel,=cpcolor,=cpdelta);
\end{euleroutput}
\begin{eulerprompt}
>fcontour("d4",xmin=-1.5,xmax=1.5,ymin=-1.5,ymax=1.5,hue=1);
\end{eulerprompt}
\begin{euleroutput}
  Illegal function result in map.
   %ploteval2:
      return map(f$,x,y;args());
  Try "trace errors" to inspect local variables after errors.
  fcontour:
      Z=%ploteval2(f$,X,Y,maps;args());
\end{euleroutput}
\begin{eulerprompt}
>P=(A_B_C_D)'; plot2d(P[1],P[2],points=1,add=1,color=12);
>insimg;
\end{eulerprompt}
\begin{eulercomment}
Masih ada minimum dan dicapai bukan di titik A, B, C, ataupun D:
\end{eulercomment}
\begin{eulerprompt}
>function f(x):=d4(x[1],x[2])
>neldermin("f",[0.2,0.2])
\end{eulerprompt}
\begin{euleroutput}
  Function call to "%nelderf" failed in nelder.
  Try "trace errors" to inspect local variables after errors.
  neldermin:
      v1=_nelder("%nelderf",v,d,eps;f$,args());
\end{euleroutput}
\begin{eulercomment}
Tampaknya pada kasus ini, koordinat titik optimalnya rasional atau hampir rasional...

Sekarang jika ABCD adalah persegi, kita harapkan titik optimalnya adalah pusat ABCD:
\end{eulercomment}
\begin{eulerprompt}
>C=[-1,1];
>plot3d("d4",xmin=-1,xmax=1,ymin=-1,ymax=1):
\end{eulerprompt}
\begin{euleroutput}
  Illegal function result in map.
   %ploteval2:
      return map(f$,x,y;args());
  f3dplot:
      z=%ploteval2(f$,x,y,maps;args());
  Try "trace errors" to inspect local variables after errors.
  plot3d:
      =cp,=cplevel,=cpcolor,=cpdelta);
\end{euleroutput}
\begin{eulerprompt}
>fcontour("d4",xmin=-1.5,xmax=1.5,ymin=-1.5,ymax=1.5,hue=1);
\end{eulerprompt}
\begin{euleroutput}
  Illegal function result in map.
   %ploteval2:
      return map(f$,x,y;args());
  Try "trace errors" to inspect local variables after errors.
  fcontour:
      Z=%ploteval2(f$,X,Y,maps;args());
\end{euleroutput}
\begin{eulerprompt}
>P=(A_B_C_D)'; plot2d(P[1],P[2],add=1,color=12,points=1);
>insimg;
\end{eulerprompt}
\begin{eulercomment}
Hasil yang diharapkan: untuk empat titik A=[-1,-1], B=[1,-1], C=[-0.5,1], D=[1,1],
neldermin mendapatkan titik sekitar [0.142858, 0.142857] (kira-kira [1/7,1/7]).
Untuk persegi (C=[-1,1]) titik optimalnya adalah pusat persegi, [0,0].

\begin{eulercomment}
\eulerheading{Contoh 7: Bola Dandelin dengan Povray}
\begin{eulercomment}
\end{eulercomment}
\eulersubheading{Uraian Materi}
\begin{eulercomment}
Teorema Dandelin: bila sebuah bidang memotong kerucut tegak membentuk elips, maka terdapat dua
bola yang menyinggung kerucut dan bidang tersebut. Titik-titik singgung kedua bola dengan
bidang adalah titik fokus elips. Untuk titik T pada elips, TF1 + TF2 sama dengan jarak antara
dua lingkaran singgung bola dengan kerucut, diukur sepanjang garis pelukis, sehingga konstan.
Ini menjelaskan mengapa irisan kerucut ini adalah elips (jumlah jarak ke dua fokus konstan).

\end{eulercomment}
\eulersubheading{Penyelesaian dengan EMT}
\begin{eulercomment}
Anda dapat menjalankan demonstrasi ini, jika Anda telah memasang Povray, dan pvengine.exe berada
di path program.

Pertama kita hitung jari-jari bola-bola tersebut.

Jika Anda melihat gambar di bawah, Anda melihat bahwa kita membutuhkan dua lingkaran yang
menyinggung dua garis yang membentuk kerucut, dan satu garis yang membentuk bidang yang
memotong kerucut.

Kita pakai file geometry.e dari Euler untuk ini.
\end{eulercomment}
\begin{eulerprompt}
>load geometry;
\end{eulerprompt}
\begin{eulercomment}
Pertama dua garis yang membentuk kerucut.
\end{eulercomment}
\begin{eulerprompt}
>g1 &= lineThrough([0,0],[1,a])
\end{eulerprompt}
\begin{euleroutput}
  
                               [- a, 1, 0]
  
\end{euleroutput}
\begin{eulerprompt}
>g2 &= lineThrough([0,0],[-1,a])
\end{eulerprompt}
\begin{euleroutput}
  
                              [- a, - 1, 0]
  
\end{euleroutput}
\begin{eulercomment}
Kemudian garis ketiga.
\end{eulercomment}
\begin{eulerprompt}
>g &= lineThrough([-1,0],[1,1])
\end{eulerprompt}
\begin{euleroutput}
  
                               [- 1, 2, 1]
  
\end{euleroutput}
\begin{eulercomment}
Kita gambar semuanya sejauh ini.
\end{eulercomment}
\begin{eulerprompt}
>setPlotRange(-1,1,0,2);
>color(black); plotLine(g(),"")
>a:=2; color(blue); plotLine(g1(),""), plotLine(g2(),""):
\end{eulerprompt}
\begin{eulercomment}
Sekarang kita ambil sebuah titik umum pada sumbu-y.
\end{eulercomment}
\begin{eulerprompt}
>P &= [0,u]
\end{eulerprompt}
\begin{euleroutput}
  
                                  [0, u]
  
\end{euleroutput}
\begin{eulercomment}
Hitung jaraknya ke g1.
\end{eulercomment}
\begin{eulerprompt}
>d1 &= distance(P,projectToLine(P,g1)); $d1
\end{eulerprompt}
\begin{eulercomment}
Hitung jaraknya ke g.
\end{eulercomment}
\begin{eulerprompt}
>d &= distance(P,projectToLine(P,g)); $d
\end{eulerprompt}
\begin{eulercomment}
Dan carilah pusat kedua lingkaran, yang jaraknya sama.
\end{eulercomment}
\begin{eulerprompt}
>sol &= solve(d1^2=d^2,u); $sol
\end{eulerprompt}
\begin{eulercomment}
Ada dua solusi.

Kita evaluasi solusi simbolik tersebut, dan temukan kedua pusat, dan kedua jarak.
\end{eulercomment}
\begin{eulerprompt}
>u := sol()
\end{eulerprompt}
\begin{euleroutput}
  [0.333333,  1]
\end{euleroutput}
\begin{eulerprompt}
>dd := d()
\end{eulerprompt}
\begin{euleroutput}
  [0.149071,  0.447214]
\end{euleroutput}
\begin{eulercomment}
Gambar lingkaran-lingkaran ke dalam gambar.
\end{eulercomment}
\begin{eulerprompt}
>color(red);
>plotCircle(circleWithCenter([0,u[1]],dd[1]),"");
>plotCircle(circleWithCenter([0,u[2]],dd[2]),"");
>insimg;
\end{eulerprompt}
\begin{eulercomment}
Hasil yang diharapkan: u = [0.333333, 1]; jari-jari dd = [0.149071, 0.447214].

\end{eulercomment}
\eulersubheading{Gambar dengan Povray}
\begin{eulercomment}
Selanjutnya kita gambar semuanya dengan Povray. Perhatikan bahwa Anda dapat mengubah perintah
apa pun dalam rangkaian perintah Povray berikut, dan menjalankan ulang semua perintah dengan
Shift-Return.

Pertama kita muat fungsi-fungsi povray.
\end{eulercomment}
\begin{eulerprompt}
>load povray;
>defaultpovray="C:\(\backslash\)Program Files\(\backslash\)POV-Ray\(\backslash\)v3.7\(\backslash\)bin\(\backslash\)pvengine.exe"
\end{eulerprompt}
\begin{euleroutput}
  C:\(\backslash\)Program Files\(\backslash\)POV-Ray\(\backslash\)v3.7\(\backslash\)bin\(\backslash\)pvengine.exe
\end{euleroutput}
\begin{eulercomment}
Kita siapkan adegan (scene) dengan tepat.
\end{eulercomment}
\begin{eulerprompt}
>povstart(zoom=11,center=[0,0,0.5],height=10°,angle=140°);
\end{eulerprompt}
\begin{euleroutput}
  Variable  not found!
  Error in degree
  Error in:
  povstart(zoom=11,center=[0,0,0.5],height=10°,angle=140°); ...
                                               ^
\end{euleroutput}
\begin{eulercomment}
Selanjutnya kita tulis kedua bola ke file Povray.
\end{eulercomment}
\begin{eulerprompt}
>writeln(povsphere([0,0,u[1]],dd[1],povlook(red)));
\end{eulerprompt}
\begin{euleroutput}
  sphere \{ <0,0,0.333333>, 0.1490711985 texture \{ pigment \{ color rgb <0.564706,0.0627451,0.0627451> \}  \} 
   finish \{ ambient 0.2 \} 
  
  \}
\end{euleroutput}
\begin{eulerprompt}
>writeln(povsphere([0,0,u[2]],dd[2],povlook(red)));
\end{eulerprompt}
\begin{euleroutput}
  sphere \{ <0,0,1>, 0.4472135955 texture \{ pigment \{ color rgb <0.564706,0.0627451,0.0627451> \}  \} 
   finish \{ ambient 0.2 \} 
  
  \}
\end{euleroutput}
\begin{eulercomment}
Dan kerucutnya, transparan.
\end{eulercomment}
\begin{eulerprompt}
>writeln(povcone([0,0,0],0,[0,0,a],1,povlook(lightgray,1)));
\end{eulerprompt}
\begin{euleroutput}
  cone \{ <0,0,0>, 0 <0,0,2>, 1 texture \{ pigment \{ color rgbf <0.705882,0.705882,0.705882,1> \}  \} 
   finish \{ ambient 0.2 \} 
   \}
\end{euleroutput}
\begin{eulercomment}
Kita membuat bidang yang dibatasi pada kerucut.
\end{eulercomment}
\begin{eulerprompt}
>gp=g();
>pc=povcone([0,0,0],0,[0,0,a],1,"");
>vp=[gp[1],0,gp[2]]; dp=gp[3];
>writeln(povplane(vp,dp,povlook(blue,0.5),pc));
\end{eulerprompt}
\begin{euleroutput}
  plane \{ <-1,0,2>, 0.4472135955
  clipped_by \{ cone \{ <0,0,0>, 0 <0,0,2>, 1 \} \}  texture \{ pigment \{ color rgbf <0.0627451,0.0627451,0.564706,0.5> \}  \} 
   finish \{ ambient 0.2 \} 
  
  \}
\end{euleroutput}
\begin{eulercomment}
Sekarang kita buat dua titik pada lingkaran-lingkaran, di mana bola-bola menyinggung kerucut.
\end{eulercomment}
\begin{eulerprompt}
>function turnz(v) := return [-v[2],v[1],v[3]]
>P1=projectToLine([0,u[1]],g1()); P1=turnz([P1[1],0,P1[2]]);
>writeln(povpoint(P1,povlook(yellow)));
\end{eulerprompt}
\begin{euleroutput}
  sphere \{ <0,0.133333,0.266667>, 0.02 texture \{ pigment \{ color rgb <0.941176,0.941176,0.392157> \}  \} 
   finish \{ ambient 0.2 \} 
  
  \}
\end{euleroutput}
\begin{eulerprompt}
>P2=projectToLine([0,u[2]],g1()); P2=turnz([P2[1],0,P2[2]]);
>writeln(povpoint(P2,povlook(yellow)));
\end{eulerprompt}
\begin{euleroutput}
  sphere \{ <0,0.4,0.8>, 0.02 texture \{ pigment \{ color rgb <0.941176,0.941176,0.392157> \}  \} 
   finish \{ ambient 0.2 \} 
  
  \}
\end{euleroutput}
\begin{eulercomment}
Kemudian kita buat dua titik di mana bola-bola menyinggung bidang.
Titik-titik ini adalah fokus elips.
\end{eulercomment}
\begin{eulerprompt}
>P3=projectToLine([0,u[1]],g()); P3=[P3[1],0,P3[2]];
>writeln(povpoint(P3,povlook(yellow)));
\end{eulerprompt}
\begin{euleroutput}
  sphere \{ <-0.0666667,0,0.466667>, 0.02 texture \{ pigment \{ color rgb <0.941176,0.941176,0.392157> \}  \} 
   finish \{ ambient 0.2 \} 
  
  \}
\end{euleroutput}
\begin{eulerprompt}
>P4=projectToLine([0,u[2]],g()); P4=[P4[1],0,P4[2]];
>writeln(povpoint(P4,povlook(yellow)));
\end{eulerprompt}
\begin{euleroutput}
  sphere \{ <0.2,0,0.6>, 0.02 texture \{ pigment \{ color rgb <0.941176,0.941176,0.392157> \}  \} 
   finish \{ ambient 0.2 \} 
  
  \}
\end{euleroutput}
\begin{eulercomment}
Berikutnya kita hitung perpotongan P1P2 dengan bidang.
\end{eulercomment}
\begin{eulerprompt}
>t1=scalp(vp,P1)-dp; t2=scalp(vp,P2)-dp; P5=P1+t1/(t1-t2)*(P2-P1);
>writeln(povpoint(P5,povlook(yellow)));
\end{eulerprompt}
\begin{euleroutput}
  sphere \{ <0,0.25,0.5>, 0.02 texture \{ pigment \{ color rgb <0.941176,0.941176,0.392157> \}  \} 
   finish \{ ambient 0.2 \} 
  
  \}
\end{euleroutput}
\begin{eulercomment}
Kita hubungkan titik-titik dengan ruas garis.
\end{eulercomment}
\begin{eulerprompt}
>writeln(povsegment(P1,P2,povlook(yellow)));
\end{eulerprompt}
\begin{euleroutput}
  cylinder \{ <0,0.133333,0.266667>, <0,0.4,0.8>, 0.01
   texture \{ pigment \{ color rgb <0.941176,0.941176,0.392157> \}  \} 
   finish \{ ambient 0.2 \} 
   \}
\end{euleroutput}
\begin{eulerprompt}
>writeln(povsegment(P5,P3,povlook(yellow)));
\end{eulerprompt}
\begin{euleroutput}
  cylinder \{ <0,0.25,0.5>, <-0.0666667,0,0.466667>, 0.01
   texture \{ pigment \{ color rgb <0.941176,0.941176,0.392157> \}  \} 
   finish \{ ambient 0.2 \} 
   \}
\end{euleroutput}
\begin{eulerprompt}
>writeln(povsegment(P5,P4,povlook(yellow)));
\end{eulerprompt}
\begin{euleroutput}
  cylinder \{ <0,0.25,0.5>, <0.2,0,0.6>, 0.01
   texture \{ pigment \{ color rgb <0.941176,0.941176,0.392157> \}  \} 
   finish \{ ambient 0.2 \} 
   \}
\end{euleroutput}
\begin{eulercomment}
Sekarang kita buat pita abu-abu, di mana bola-bola menyinggung kerucut.
\end{eulercomment}
\begin{eulerprompt}
>pcw=povcone([0,0,0],0,[0,0,a],1.01);
>pc1=povcylinder([0,0,P1[3]-defaultpointsize/2],[0,0,P1[3]+defaultpointsize/2],1);
>writeln(povintersection([pcw,pc1],povlook(gray)));
\end{eulerprompt}
\begin{euleroutput}
  intersection \{ cone \{ <0,0,0>, 0 <0,0,2>, 1.01 texture \{ pigment \{ color rgb <0.470588,0.470588,0.470588> \}  \} 
   finish \{ ambient 0.2 \} 
   \}
  cylinder \{ <0,0,0.256667>, <0,0,0.276667>, 1
   texture \{ pigment \{ color rgb <0.470588,0.470588,0.470588> \}  \} 
   finish \{ ambient 0.2 \} 
   \}
   texture \{ pigment \{ color rgb <0.470588,0.470588,0.470588> \}  \} 
   finish \{ ambient 0.2 \} 
   \}
\end{euleroutput}
\begin{eulerprompt}
>pc2=povcylinder([0,0,P2[3]-defaultpointsize/2],[0,0,P2[3]+defaultpointsize/2],1);
>writeln(povintersection([pcw,pc2],povlook(gray)));
\end{eulerprompt}
\begin{euleroutput}
  intersection \{ cone \{ <0,0,0>, 0 <0,0,2>, 1.01 texture \{ pigment \{ color rgb <0.470588,0.470588,0.470588> \}  \} 
   finish \{ ambient 0.2 \} 
   \}
  cylinder \{ <0,0,0.79>, <0,0,0.81>, 1
   texture \{ pigment \{ color rgb <0.470588,0.470588,0.470588> \}  \} 
   finish \{ ambient 0.2 \} 
   \}
   texture \{ pigment \{ color rgb <0.470588,0.470588,0.470588> \}  \} 
   finish \{ ambient 0.2 \} 
   \}
\end{euleroutput}
\begin{eulercomment}
Jalankan program Povray.
\end{eulercomment}
\begin{eulerprompt}
>povend();
\end{eulerprompt}
\begin{euleroutput}
  exec:
      return _exec(program,param,dir,print,hidden,wait);
  povray:
      exec(program,params,defaulthome);
  Try "trace errors" to inspect local variables after errors.
  povend:
      povray(file,w,h,aspect,exit); 
\end{euleroutput}
\begin{eulercomment}
Untuk mendapatkan anaglif, kita perlu memasukkan semuanya ke dalam sebuah fungsi adegan.
Fungsi ini akan dipakai dua kali nanti.
\end{eulercomment}
\begin{eulerprompt}
>function scene () ...
\end{eulerprompt}
\begin{eulerudf}
  global a,u,dd,g,g1,defaultpointsize;
  writeln(povsphere([0,0,u[1]],dd[1],povlook(red)));
  writeln(povsphere([0,0,u[2]],dd[2],povlook(red)));
  writeln(povcone([0,0,0],0,[0,0,a],1,povlook(lightgray,1)));
  gp=g();
  pc=povcone([0,0,0],0,[0,0,a],1,"");
  vp=[gp[1],0,gp[2]]; dp=gp[3];
  writeln(povplane(vp,dp,povlook(blue,0.5),pc));
  P1=projectToLine([0,u[1]],g1()); P1=turnz([P1[1],0,P1[2]]);
  writeln(povpoint(P1,povlook(yellow)));
  P2=projectToLine([0,u[2]],g1()); P2=turnz([P2[1],0,P2[2]]);
  writeln(povpoint(P2,povlook(yellow)));
  P3=projectToLine([0,u[1]],g()); P3=[P3[1],0,P3[2]];
  writeln(povpoint(P3,povlook(yellow)));
  P4=projectToLine([0,u[2]],g()); P4=[P4[1],0,P4[2]];
  writeln(povpoint(P4,povlook(yellow)));
  t1=scalp(vp,P1)-dp; t2=scalp(vp,P2)-dp; P5=P1+t1/(t1-t2)*(P2-P1);
  writeln(povpoint(P5,povlook(yellow)));
  writeln(povsegment(P1,P2,povlook(yellow)));
  writeln(povsegment(P5,P3,povlook(yellow)));
  writeln(povsegment(P5,P4,povlook(yellow)));
  pcw=povcone([0,0,0],0,[0,0,a],1.01);
  pc1=povcylinder([0,0,P1[3]-defaultpointsize/2],[0,0,P1[3]+defaultpointsize/2],1);
  writeln(povintersection([pcw,pc1],povlook(gray)));
  pc2=povcylinder([0,0,P2[3]-defaultpointsize/2],[0,0,P2[3]+defaultpointsize/2],1);
  writeln(povintersection([pcw,pc2],povlook(gray)));
  endfunction
\end{eulerudf}
\begin{eulercomment}
Anda memerlukan kacamata merah/cyan untuk menikmati efek berikut.
\end{eulercomment}
\begin{eulerprompt}
>povanaglyph("scene",zoom=11,center=[0,0,0.5],height=10°,angle=140°);
\end{eulerprompt}
\begin{euleroutput}
  Variable  not found!
  Error in degree
  Error in:
  ... lyph("scene",zoom=11,center=[0,0,0.5],height=10°,angle=140°) ...
                                                       ^
\end{euleroutput}
\begin{eulercomment}
Catatan: bagian Povray hanya berjalan bila POV-Ray terpasang dan jalurnya benar. Jika muncul
pesan "Command was not allowed!" atau galat exec, berarti Povray tidak dapat dijalankan dari
EMT pada komputer Anda; bagian gambar 2D di atas tetap berfungsi.

\begin{eulercomment}
\eulerheading{Contoh 8: Geometri Bumi}
\begin{eulercomment}
\end{eulercomment}
\eulersubheading{Uraian Materi}
\begin{eulercomment}
Pada bola dengan jari-jari R, jarak terpendek dua titik adalah busur lingkaran besar (great
circle), yaitu lingkaran dengan pusat di pusat bola. Untuk titik (lintang f1, bujur l1) dan
(f2, l2), sudut pusat d memenuhi hukum kosinus bola:
\end{eulercomment}
\begin{eulerformula}
\[
\cos d = \sin\varphi_1\sin\varphi_2+\cos\varphi_1\cos\varphi_2\cos(\lambda_2-\lambda_1),
\]
\end{eulerformula}
\begin{eulercomment}
dan jarak = R.d. Lintasan lingkaran besar tidak mempunyai arah (heading) konstan.
Jumlah sudut segitiga bola lebih dari 180°; selisihnya disebut ekses sferis E = A+B+C-pi, dan
luas segitiga bola = R\textasciicircum{}2 . E. Bumi sebenarnya berbentuk elipsoid, sehingga fungsi dengan awalan
"es" (esdist, esdir, esadd, esarea) memakai hampiran yang lebih baik daripada bola ideal.

\end{eulercomment}
\eulersubheading{Penyelesaian dengan EMT}
\begin{eulercomment}
Pada buku catatan ini, kita akan melakukan beberapa perhitungan bola. Fungsi-fungsinya
terdapat dalam file "spherical.e" pada folder examples. Kita perlu memuat file tersebut
terlebih dahulu.
\end{eulercomment}
\begin{eulerprompt}
>load "spherical.e";
\end{eulerprompt}
\begin{eulercomment}
Untuk memasukkan posisi geografis, kita pakai vektor dengan dua koordinat dalam radian (utara
dan timur, nilai negatif untuk selatan dan barat). Berikut koordinat Kampus FMIPA UNY.
\end{eulercomment}
\begin{eulerprompt}
>FMIPA=[rad(-7,-46.467),rad(110,23.05)]
\end{eulerprompt}
\begin{euleroutput}
  [-0.13569,  1.92657]
\end{euleroutput}
\begin{eulercomment}
Anda dapat mencetak posisi ini dengan sposprint (spherical position print).
\end{eulercomment}
\begin{eulerprompt}
>sposprint(FMIPA) // posisi garis lintang dan garis bujur FMIPA UNY
\end{eulerprompt}
\begin{euleroutput}
  S 7°46.467' E 110°23.050'
\end{euleroutput}
\begin{eulercomment}
Mari kita tambahkan dua kota lagi, Solo dan Semarang.
\end{eulercomment}
\begin{eulerprompt}
>Solo=[rad(-7,-34.333),rad(110,49.683)]; Semarang=[rad(-6,-59.05),rad(110,24.533)];
>sposprint(Solo), sposprint(Semarang),
\end{eulerprompt}
\begin{euleroutput}
  S 7°34.333' E 110°49.683'
  S 6°59.050' E 110°24.533'
\end{euleroutput}
\begin{eulercomment}
Pertama kita hitung vektor dari satu kota ke kota lain pada bola ideal. Vektor ini adalah
[heading,jarak] dalam radian. Untuk menghitung jarak di bumi, kita kalikan dengan jari-jari
bumi pada lintang 7°.
\end{eulercomment}
\begin{eulerprompt}
>br=svector(FMIPA,Solo); degprint(br[1]), br[2]*rearth(7°)->km // perkiraan jarak FMIPA-Solo
\end{eulerprompt}
\begin{euleroutput}
  65°20'26.60''
  Variable  not found!
  Error in degree
  Error in:
  ... or(FMIPA,Solo); degprint(br[1]), br[2]*rearth(7°)->km // perk ...
                                                       ^
\end{euleroutput}
\begin{eulercomment}
Ini adalah hampiran yang baik. Rutin-rutin berikut memakai hampiran yang lebih baik lagi.
Pada jarak yang pendek hasilnya hampir sama.

Catatan: pada versi asli terjadi galat "Commands must be separated by semicolon" bila ada
komentar "//" setelah -\textgreater{}" km". Karena itu komentar untuk perintah semacam ini ditaruh di baris
tersendiri.
Perkiraan jarak FMIPA-Semarang:
\end{eulercomment}
\begin{eulerprompt}
>esdist(FMIPA,Semarang)->" km"
\end{eulerprompt}
\begin{euleroutput}
  88.0114026318 km
\end{euleroutput}
\begin{eulercomment}
Ada fungsi untuk heading, dengan memperhitungkan bentuk elips bumi. Sekali lagi, kita cetak
dengan cara yang lebih baik.
\end{eulercomment}
\begin{eulerprompt}
>sdegprint(esdir(FMIPA,Solo))
\end{eulerprompt}
\begin{euleroutput}
       65.34°
\end{euleroutput}
\begin{eulercomment}
Sudut suatu segitiga melebihi 180° pada bola.
\end{eulercomment}
\begin{eulerprompt}
>asum=sangle(Solo,FMIPA,Semarang)+sangle(FMIPA,Solo,Semarang)+sangle(FMIPA,Semarang,Solo); degprint(asum)
\end{eulerprompt}
\begin{euleroutput}
  180°0'10.77''
\end{euleroutput}
\begin{eulercomment}
Ini dapat dipakai untuk menghitung luas segitiga. Catatan: untuk segitiga kecil, ini tidak
akurat karena galat pengurangan pada asum-pi.
Perkiraan luas segitiga FMIPA-Solo-Semarang:
\end{eulercomment}
\begin{eulerprompt}
>(asum-pi)*rearth(48°)^2->" km^2"
\end{eulerprompt}
\begin{euleroutput}
  Variable  not found!
  Error in degree
  Error in:
  (asum-pi)*rearth(48°)^2->" km^2" ...
                       ^
\end{euleroutput}
\begin{eulercomment}
Ada fungsi untuk ini, yang memakai lintang rata-rata segitiga untuk menghitung jari-jari bumi,
dan menangani galat pembulatan untuk segitiga yang sangat kecil.
Perkiraan yang sama dengan fungsi esarea():
\end{eulercomment}
\begin{eulerprompt}
>esarea(Solo,FMIPA,Semarang)->" km^2"
\end{eulerprompt}
\begin{euleroutput}
  2123.64310526 km^2
\end{euleroutput}
\begin{eulercomment}
Kita juga dapat menambahkan vektor ke posisi. Vektor memuat heading dan jarak, keduanya dalam
radian. Untuk mendapatkan vektor, kita pakai svector. Untuk menambahkan vektor ke posisi, kita
pakai saddvector.
\end{eulercomment}
\begin{eulerprompt}
>v=svector(FMIPA,Solo); sposprint(saddvector(FMIPA,v)), sposprint(Solo),
\end{eulerprompt}
\begin{euleroutput}
  S 7°34.333' E 110°49.683'
  S 7°34.333' E 110°49.683'
\end{euleroutput}
\begin{eulercomment}
Fungsi-fungsi ini mengandaikan bola ideal. Hal yang sama di bumi.
\end{eulercomment}
\begin{eulerprompt}
>sposprint(esadd(FMIPA,esdir(FMIPA,Solo),esdist(FMIPA,Solo))), sposprint(Solo),
\end{eulerprompt}
\begin{euleroutput}
  S 7°34.333' E 110°49.683'
  S 7°34.333' E 110°49.683'
\end{euleroutput}
\begin{eulercomment}
Mari kita beralih ke contoh yang lebih besar, Tugu Jogja dan Monas Jakarta (menggunakan Google
Earth untuk mencari koordinatnya).
\end{eulercomment}
\begin{eulerprompt}
>Tugu=[-7.7833°,110.3661°]; Monas=[-6.175°,106.811944°];
\end{eulerprompt}
\begin{euleroutput}
  Variable  not found!
  Error in degree
  Error in:
  Tugu=[-7.7833°,110.3661°]; Monas=[-6.175°,106.811944°]; ...
                 ^
\end{euleroutput}
\begin{eulerprompt}
>sposprint(Tugu), sposprint(Monas)
\end{eulerprompt}
\begin{euleroutput}
  Variable or function Tugu not found.
  Error in:
  sposprint(Tugu), sposprint(Monas) ...
                ^
\end{euleroutput}
\begin{eulercomment}
Menurut Google Earth, jaraknya 429.66km. Kita mendapatkan hampiran yang baik.
Perkiraan jarak Tugu Jogja - Monas Jakarta:
\end{eulercomment}
\begin{eulerprompt}
>esdist(Tugu,Monas)->" km"
\end{eulerprompt}
\begin{euleroutput}
  Variable or function Tugu not found.
  Error in:
  esdist(Tugu,Monas)->" km" ...
             ^
\end{euleroutput}
\begin{eulercomment}
Heading-nya sama dengan yang dihitung di Google Earth.
\end{eulercomment}
\begin{eulerprompt}
>degprint(esdir(Tugu,Monas))
\end{eulerprompt}
\begin{euleroutput}
  Variable or function Tugu not found.
  Error in:
  degprint(esdir(Tugu,Monas)) ...
                     ^
\end{euleroutput}
\begin{eulercomment}
Namun, kita tidak lagi mendapatkan posisi tujuan yang tepat, jika kita menambahkan heading dan
jarak ke posisi semula. Ini karena kita tidak menghitung fungsi inversnya secara tepat, melainkan
mengambil hampiran jari-jari bumi sepanjang lintasan.
\end{eulercomment}
\begin{eulerprompt}
>sposprint(esadd(Tugu,esdir(Tugu,Monas),esdist(Tugu,Monas)))
\end{eulerprompt}
\begin{euleroutput}
  Variable or function Tugu not found.
  Error in:
  sposprint(esadd(Tugu,esdir(Tugu,Monas),esdist(Tugu,Monas))) ...
                      ^
\end{euleroutput}
\begin{eulercomment}
Galatnya tidak besar.
\end{eulercomment}
\begin{eulerprompt}
>sposprint(Monas),
\end{eulerprompt}
\begin{euleroutput}
  Variable or function Monas not found.
  Error in:
  sposprint(Monas), ...
                 ^
\end{euleroutput}
\begin{eulercomment}
Tentu saja, kita tidak bisa berlayar dengan heading yang sama dari satu tujuan ke tujuan lain,
jika kita ingin mengambil lintasan terpendek. Bayangkan, Anda terbang ke arah timur laut
mulai dari sebarang titik di bumi. Maka Anda akan berputar spiral menuju kutub utara.
Lingkaran besar tidak mengikuti heading yang konstan!

Perhitungan berikut menunjukkan bahwa kita jauh meleset dari tujuan yang benar, jika kita
memakai heading yang sama selama perjalanan.
\end{eulercomment}
\begin{eulerprompt}
>dist=esdist(Tugu,Monas); hd=esdir(Tugu,Monas);
\end{eulerprompt}
\begin{euleroutput}
  Variable or function Tugu not found.
  Error in:
  dist=esdist(Tugu,Monas); hd=esdir(Tugu,Monas); ...
                  ^
\end{euleroutput}
\begin{eulercomment}
Sekarang kita tambahkan 10 kali sepersepuluh jarak, memakai heading ke Monas yang kita
dapatkan di Tugu.
\end{eulercomment}
\begin{eulerprompt}
>p=Tugu; loop 1 to 10; p=esadd(p,hd,dist/10); end;
\end{eulerprompt}
\begin{euleroutput}
  Variable Tugu not found!
  Error in:
  p=Tugu; loop 1 to 10; p=esadd(p,hd,dist/10); end; ...
        ^
\end{euleroutput}
\begin{eulercomment}
Hasilnya jauh meleset.
\end{eulercomment}
\begin{eulerprompt}
>sposprint(p), skmprint(esdist(p,Monas))
\end{eulerprompt}
\begin{euleroutput}
  Function sposprint needs a vector for pos
  Error in:
  sposprint(p), skmprint(esdist(p,Monas)) ...
              ^
\end{euleroutput}
\begin{eulercomment}
Sebagai contoh lain, mari kita ambil dua titik di bumi pada lintang yang sama.
\end{eulercomment}
\begin{eulerprompt}
>P1=[30°,10°]; P2=[30°,50°];
\end{eulerprompt}
\begin{euleroutput}
  Variable  not found!
  Error in degree
  Error in:
  P1=[30°,10°]; P2=[30°,50°]; ...
          ^
\end{euleroutput}
\begin{eulercomment}
Lintasan terpendek dari P1 ke P2 bukanlah lingkaran lintang 30°, tetapi lintasan yang lebih
pendek yang dimulai 10° lebih ke utara di P1.
\end{eulercomment}
\begin{eulerprompt}
>sdegprint(esdir(P1,P2))
\end{eulerprompt}
\begin{euleroutput}
       90.00°
\end{euleroutput}
\begin{eulercomment}
Tetapi, jika kita mengikuti pembacaan kompas ini, kita akan berputar spiral ke kutub utara!
Jadi kita harus menyesuaikan heading di sepanjang jalan. Untuk keperluan kasar, kita sesuaikan
pada setiap 1/10 dari jarak total.
\end{eulercomment}
\begin{eulerprompt}
>p=P1;  dist=esdist(P1,P2); ...
>  loop 1 to 10; dir=esdir(p,P2); sdegprint(dir), p=esadd(p,dir,dist/10); end;
\end{eulerprompt}
\begin{euleroutput}
       90.00°
       90.00°
       90.00°
       90.00°
       90.00°
       90.00°
       90.00°
       90.00°
       90.00°
       90.00°
\end{euleroutput}
\begin{eulercomment}
Jaraknya tidak tepat, karena kita menambahkan sedikit galat, jika kita mengikuti heading yang
sama terlalu lama.
\end{eulercomment}
\begin{eulerprompt}
>skmprint(esdist(p,P2))
\end{eulerprompt}
\begin{euleroutput}
       0.000km
\end{euleroutput}
\begin{eulercomment}
Kita mendapatkan hampiran yang baik, jika kita menyesuaikan heading setiap 1/100 dari jarak
total dari Tugu ke Monas.
\end{eulercomment}
\begin{eulerprompt}
>p=Tugu; dist=esdist(Tugu,Monas); ...
>  loop 1 to 100; p=esadd(p,esdir(p,Monas),dist/100); end;
\end{eulerprompt}
\begin{euleroutput}
  Variable Tugu not found!
  Error in:
  p=Tugu; dist=esdist(Tugu,Monas);   loop 1 to 100; p=esadd(p,es ...
        ^
\end{euleroutput}
\begin{eulerprompt}
>skmprint(esdist(p,Monas))
\end{eulerprompt}
\begin{euleroutput}
  Variable or function Monas not found.
  Error in:
  skmprint(esdist(p,Monas)) ...
                         ^
\end{euleroutput}
\begin{eulercomment}
Untuk keperluan navigasi, kita dapat memperoleh rangkaian posisi GPS sepanjang lingkaran besar
menuju Monas dengan fungsi navigate.
\end{eulercomment}
\begin{eulerprompt}
>load spherical; v=navigate(Tugu,Monas,10); ...
>  loop 1 to rows(v); sposprint(v[#]), end;
\end{eulerprompt}
\begin{euleroutput}
  Variable or function Tugu not found.
  Error in:
  load spherical; v=navigate(Tugu,Monas,10);   loop 1 to rows(v) ...
                                 ^
\end{euleroutput}
\begin{eulercomment}
Kita tulis sebuah fungsi, yang menggambar bumi, kedua posisi, dan posisi-posisi di antaranya.
\end{eulercomment}
\begin{eulerprompt}
>function testplot ...
\end{eulerprompt}
\begin{eulerudf}
  useglobal;
  plotearth;
  plotpos(Tugu,"Tugu Jogja"); plotpos(Monas,"Monas Jakarta");
  plotposline(v);
  endfunction
\end{eulerudf}
\begin{eulercomment}
Sekarang gambar semuanya.
\end{eulercomment}
\begin{eulerprompt}
>plot3d("testplot",angle=25, height=6,>own,>user,zoom=4):
\end{eulerprompt}
\begin{euleroutput}
  Variable or function Tugu not found.
  testplot:
      plotpos(Tugu,"Tugu Jogja"); plotpos(Monas,"Monas Jakarta");
  Try "trace errors" to inspect local variables after errors.
  plot3d:
      f(args());
\end{euleroutput}
\begin{eulercomment}
Atau pakai plot3d untuk mendapatkan tampilan anaglif. Ini tampak sangat bagus dengan
kacamata merah/cyan.
\end{eulercomment}
\begin{eulerprompt}
>plot3d("testplot",angle=25,height=6,distance=5,own=1,anaglyph=1,zoom=4):
\end{eulerprompt}
\begin{euleroutput}
  Variable or function Tugu not found.
  testplot:
      plotpos(Tugu,"Tugu Jogja"); plotpos(Monas,"Monas Jakarta");
  Try "trace errors" to inspect local variables after errors.
  plot3d:
      f(args());
\end{euleroutput}
\begin{eulercomment}
Hasil yang diharapkan (bandingkan keluaran asli): jarak FMIPA-Solo = 53.8945 km, heading 65°20'26.60'';
luas segitiga FMIPA-Solo-Semarang = 2123.64 km\textasciicircum{}2; jarak Tugu-Monas mendekati 429.66 km dengan
heading 294°17'2.85''; setelah 10 langkah dengan heading tetap, posisi meleset 1.529 km.

\begin{eulercomment}
\eulerheading{Latihan (Soal Akhir)}
\begin{eulercomment}
\end{eulercomment}
\eulersubheading{Latihan 1: Segi-n beraturan jika diketahui pusat O, n, dan jari-jari lingkaran luar r}
\begin{eulercomment}
Uraian materi: segi-n beraturan mempunyai n titik sudut pada lingkaran luar berpusat O, dengan
sudut pusat yang menghadap setiap sisi sebesar 360/n derajat (2pi/n radian). Titik sudut ke-k:
\end{eulercomment}
\begin{eulerformula}
\[
V_k = O + r(\cos(\theta_0 + 2\pi k/n),\ \sin(\theta_0 + 2\pi k/n)),\quad k=0,\dots,n-1.
\]
\end{eulerformula}
\begin{eulercomment}
Sama dengan petunjuk soal: titik sudut adalah perpotongan lingkaran luar dengan garis-garis
yang melalui pusat dan saling membentuk sudut kelipatan 360/n. Untuk n ganjil, satu titik sudut
di atas (theta0 = 90°); untuk n genap, dua titik sudut di kiri dan kanan pusat (theta0 = 0°).
Panjang sisi = 2 r sin(pi/n).
\end{eulercomment}
\begin{eulerprompt}
>function titikSegiN (O,n,r) ...
\end{eulerprompt}
\begin{eulerudf}
    if mod(n,2)==1 then t0=pi/2; else t0=0; endif;
    t=t0+2*pi*(0:n-1)/n;
    return (O[1]+r*cos(t))_(O[2]+r*sin(t));
  endfunction
\end{eulerudf}
\begin{eulerprompt}
>function gambarSegiN (O,n,r) ...
\end{eulerprompt}
\begin{eulerudf}
    V=titikSegiN(O,n,r);
    setPlotRange(O[1]-1.5*r,O[1]+1.5*r,O[2]-1.5*r,O[2]+1.5*r);
    aspect(1);
    plotCircle(circleWithCenter(O,r),"");
    plotPoint(O,"O");
    for i=1 to n
      j=i+1; if j>n then j=1; endif;
      plotPoint(V[,i]');
      plotSegment(V[,i]',V[,j]',"");
      plotSegment(O,V[,i]',"");
    end;
  endfunction
\end{eulerudf}
\begin{eulerprompt}
>gambarSegiN([0,0],3,2):
>gambarSegiN([0,0],4,2):
>gambarSegiN([1,1],5,2):
>gambarSegiN([0,0],6,2):
>gambarSegiN([0,0],7,2):
>gambarSegiN([0,0],8,2):
\end{eulerprompt}
\begin{eulercomment}
Pemeriksaan panjang sisi untuk r = 2, n = 5:
\end{eulercomment}
\begin{eulerprompt}
>V=titikSegiN([0,0],5,2); distance(V[,1]',V[,2]'), 2*2*sin(pi/5)
\end{eulerprompt}
\begin{euleroutput}
  2.35114100917
  2.35114100917
\end{euleroutput}
\begin{eulercomment}
Hasil yang diharapkan: kedua nilai 2.35114. Panjang sisi untuk r=2: n=3: 3.4641; n=4: 2.8284;
n=5: 2.3511; n=6: 2 (sama dengan r); n=7: 1.7355; n=8: 1.5307.
Penjelasan: untuk n=6 sisi sama dengan jari-jari karena segitiga pusat-dua titik sudut berurutan
sama sisi (sudut pusat 60°). Makin besar n, segi-n makin mendekati lingkaran.

\end{eulercomment}
\eulersubheading{Latihan 2: Parabola yang melalui 3 titik yang diketahui}
\begin{eulercomment}
Uraian materi: parabola y = ax\textasciicircum{}2+bx+c memiliki tiga koefisien, sehingga tiga titik (dengan
absis berbeda) menentukan parabola tersebut secara tunggal. Substitusi tiap titik menghasilkan
sistem persamaan linear (SPL) 3 x 3 untuk a, b, c.

Contoh: titik-titik (-2,5), (0,-1), (3,2). SPL: 4a-2b+c=5, c=-1, 9a+3b+c=2.
\end{eulercomment}
\begin{eulerprompt}
>xp=[-2,0,3]; yp=[5,-1,2];
>M=(xp^2)'|xp'|ones(3,1)
\end{eulerprompt}
\begin{euleroutput}
              4            -2             1 
              0             0             1 
              9             3             1 
\end{euleroutput}
\begin{eulerprompt}
>koef=M\(\backslash\)yp'
\end{eulerprompt}
\begin{euleroutput}
            0.8 
           -1.4 
             -1 
\end{euleroutput}
\begin{eulerprompt}
>pa=koef[1]; pb=koef[2]; pc=koef[3];
>pa*xp^2+pb*xp+pc
\end{eulerprompt}
\begin{euleroutput}
  [5,  -1,  2]
\end{euleroutput}
\begin{eulerprompt}
>plot2d("pa*x^2+pb*x+pc",-3,4); plot2d(xp,yp,points=1,add=1,color=12):
\end{eulerprompt}
\begin{eulercomment}
Hasil yang diharapkan: koef = [0.8; -1.4; -1], yaitu y = 0.8x\textasciicircum{}2 - 1.4x - 1; nilai parabola
di xp adalah [5, -1, 2], sama dengan yp.
Penyelesaian eksak dengan Maxima:
\end{eulercomment}
\begin{eulerprompt}
>&remvalue(a,b,c);
>sol &= solve([a*(-2)^2+b*(-2)+c=5, c=-1, a*3^2+b*3+c=2],[a,b,c]); $sol
>parabolanya &= (a*x^2+b*x+c) with sol[1]; $'y=parabolanya
>plot2d(&parabolanya,-3,4):
\end{eulerprompt}
\begin{eulercomment}
Hasil yang diharapkan: a=4/5, b=-7/5, c=-1.
Penjelasan: kurva melalui ketiga titik. Karena a = 0.8 \textgreater{} 0 parabola terbuka ke atas, dengan
titik puncak di x = -b/(2a) = 0.875.

\end{eulercomment}
\eulersubheading{Latihan 3: Segi-4 garis singgung}
\begin{eulercomment}
Uraian materi: segi-4 garis singgung adalah segi-4 yang keempat sisinya menyinggung satu
lingkaran (lingkaran dalam). Syarat perlu dan cukup (teorema Pitot, untuk segi-4 cembung):
jumlah panjang sisi-sisi yang berhadapan sama,
\end{eulercomment}
\begin{eulerformula}
\[
AB + CD = BC + DA.
\]
\end{eulerformula}
\begin{eulercomment}
Bukti (perlu): jika lingkaran dalam menyinggung sisi-sisi, panjang garis singgung dari satu titik
sudut ke lingkaran sama untuk kedua sisi yang bertemu di titik itu. Misal panjangnya p, q, r, s
di A, B, C, D. Maka AB = p+q, BC = q+r, CD = r+s, DA = s+p, sehingga AB+CD = p+q+r+s = BC+DA.

Catatan: pada soal tertulis "hasil kali panjang sisi-sisi yang berhadapan sama". Ini kurang tepat:
syarat yang benar adalah hasil JUMLAH. Contoh penyangkal untuk "hasil kali": sisi 3, 2, 4, 6 (secara
berurutan) memberi AB x CD = 3 x 4 = 12 = 2 x 6 = BC x DA, padahal AB+CD = 7 tidak sama dengan BC+DA = 8, jadi
segi-4 tersebut bukan segi-4 garis singgung. Selain itu, setiap layang-layang mempunyai hasil kali
yang sama dengan sendirinya (AB.CD = BC.DA) sehingga hasil kali tidak dapat membedakannya.
Karena itu pada latihan ini yang diuji adalah jumlahnya (syarat Pitot).

Fungsi pengujian: mengembalikan [AB+CD, BC+DA, selisih].
\end{eulercomment}
\begin{eulerprompt}
>function tesPitot (A,B,C,D) ...
\end{eulerprompt}
\begin{eulerudf}
    s1=distance(A,B)+distance(C,D);
    s2=distance(B,C)+distance(D,A);
    return [s1,s2,s1-s2];
  endfunction
\end{eulerudf}
\begin{eulercomment}
Contoh 1: layang-layang A(-1,0), B(0,2), C(4,0), D(0,-2).
\end{eulercomment}
\begin{eulerprompt}
>A=[-1,0]; B=[0,2]; C=[4,0]; D=[0,-2];
>tesPitot(A,B,C,D)
\end{eulerprompt}
\begin{euleroutput}
  [6.7082,  6.7082,  0]
\end{euleroutput}
\begin{eulercomment}
Hasil yang diharapkan: [6.7082, 6.7082, 0], jadi AB+CD = BC+DA (syarat dipenuhi).
Uji kedua dengan garis bagi sudut: keempat garis bagi harus bertemu di satu titik.
\end{eulercomment}
\begin{eulerprompt}
>gA4=angleBisector(D,A,B); gB4=angleBisector(A,B,C); gC4=angleBisector(B,C,D); gD4=angleBisector(C,D,A);
>setPlotRange(-3,6,-4,4); aspect(1);
>plotPoint(A,"A"); plotPoint(B,"B"); plotPoint(C,"C"); plotPoint(D,"D");
>plotSegment(A,B,""); plotSegment(B,C,""); plotSegment(C,D,""); plotSegment(D,A,"");
>color(5); plotLine(gA4); plotLine(gB4); plotLine(gC4); plotLine(gD4); color(1);
>Pd=lineIntersection(gA4,gB4)
\end{eulerprompt}
\begin{euleroutput}
  [0.666667,  0]
\end{euleroutput}
\begin{eulerprompt}
>jarakKeGaris(Pd,gC4), jarakKeGaris(Pd,gD4)
\end{eulerprompt}
\begin{euleroutput}
  0
  0
\end{euleroutput}
\begin{eulerprompt}
>rd=[jarakKeGaris(Pd,lineThrough(A,B)),jarakKeGaris(Pd,lineThrough(B,C)),jarakKeGaris(Pd,lineThrough(C,D)),jarakKeGaris(Pd,lineThrough(D,A))]
\end{eulerprompt}
\begin{euleroutput}
  [1.49071,  1.49071,  1.49071,  1.49071]
\end{euleroutput}
\begin{eulerprompt}
>plotPoint(Pd,"I"); plotCircle(circleWithCenter(Pd,rd[1]),"Lingkaran dalam"):
\end{eulerprompt}
\begin{eulercomment}
Hasil yang diharapkan: Pd = [0.666667, 0]; jarak Pd ke gC4 dan gD4 adalah 0; keempat jarak
rd sama, yaitu 1.49071 (jari-jari lingkaran dalam).
Penjelasan: keempat garis bagi sudut konkuren di Pd dan Pd berjarak sama (1.49071) ke keempat sisi,
sehingga lingkaran berpusat Pd berjari-jari 1.49071 menyinggung keempat sisi: ini segi-4 garis
singgung, sesuai teorema Pitot (selisih jumlah sisi berhadapan = 0).

Contoh 2 (bukan segi-4 garis singgung): A(0,0), B(4,0), C(5,3), D(0,2).
\end{eulercomment}
\begin{eulerprompt}
>A=[0,0]; B=[4,0]; C=[5,3]; D=[0,2];
>tesPitot(A,B,C,D)
\end{eulerprompt}
\begin{euleroutput}
  [9.09902,  5.16228,  3.93674]
\end{euleroutput}
\begin{eulerprompt}
>gA4=angleBisector(D,A,B); gB4=angleBisector(A,B,C); gC4=angleBisector(B,C,D); gD4=angleBisector(C,D,A);
>Pd=lineIntersection(gA4,gB4)
\end{eulerprompt}
\begin{euleroutput}
  [2.32456,  2.32456]
\end{euleroutput}
\begin{eulerprompt}
>jarakKeGaris(Pd,gC4), jarakKeGaris(Pd,gD4)
\end{eulerprompt}
\begin{euleroutput}
  1.26425001148
  1.72473458997
\end{euleroutput}
\begin{eulercomment}
Hasil yang diharapkan: tesPitot = [9.0990, 5.1623, 3.9367]; Pd = [2.32456, 2.32456]; jarak ke
gC4 = 1.26425 dan ke gD4 = 1.72473 (tidak nol).
Penjelasan: keempat garis bagi tidak bertemu di satu titik, dan jumlah sisi berhadapan tidak
sama, jadi segi-4 ini bukan segi-4 garis singgung. Kedua uji (Pitot dan garis bagi) konsisten.

\end{eulercomment}
\eulersubheading{Latihan 4: Elips jika diketahui kedua titik fokus P dan Q}
\begin{eulercomment}
Uraian materi: elips dengan fokus P dan Q adalah himpunan titik T dengan TP + TQ = 2a (konstan,
2a \textgreater{} PQ). Jika c = PQ/2 dan b\textasciicircum{}2 = a\textasciicircum{}2 - c\textasciicircum{}2, pada posisi baku (pusat di titik asal, fokus di
sumbu-x) persamaannya x\textasciicircum{}2/a\textasciicircum{}2 + y\textasciicircum{}2/b\textasciicircum{}2 = 1.

(a) Fokus P(-3,0), Q(3,0), jumlah jarak 10 (a=5, c=3, b=4), digambar secara implisit.
\end{eulercomment}
\begin{eulerprompt}
>eli &= sqrt((x+3)^2+y^2)+sqrt((x-3)^2+y^2)-10; $eli='0
>P=[-3,0]; Q=[3,0];
>setPlotRange(7); aspect(1); plotPoint(P,"P"); plotPoint(Q,"Q");
>plot2d(eli,level=0,add=1,contourcolor=blue):
\end{eulerprompt}
\begin{eulercomment}
Hasil yang diharapkan: elips x\textasciicircum{}2/25 + y\textasciicircum{}2/16 = 1 dengan fokus P dan Q.
Pemeriksaan numerik dengan titik parametrik T(5 cos t, 4 sin t):
\end{eulercomment}
\begin{eulerprompt}
>t=1; T=[5*cos(t),4*sin(t)]; distance(T,P)+distance(T,Q)
\end{eulerprompt}
\begin{euleroutput}
  10
\end{euleroutput}
\begin{eulercomment}
Hasil yang diharapkan: 10.

(b) Fokus sebarang (miring): P(1,1), Q(5,4), jumlah jarak 8. Digambar secara parametrik
dengan memutar elips baku sejauh arah vektor PQ.
\end{eulercomment}
\begin{eulerprompt}
>P=[1,1]; Q=[5,4]; jj=8;
>M0=(P+Q)/2; c2=distance(P,Q)/2; a2=jj/2; b2=sqrt(a2^2-c2^2); th=atan2(Q[2]-P[2],Q[1]-P[1]);
>tt=linspace(0,2*pi,200);
>xe=M0[1]+a2*cos(tt)*cos(th)-b2*sin(tt)*sin(th); ye=M0[2]+a2*cos(tt)*sin(th)+b2*sin(tt)*cos(th);
>setPlotRange(-2,9,-3,8); aspect(1); plotPoint(P,"P"); plotPoint(Q,"Q");
>plot2d(xe,ye,add=1,color=blue):
>jum=sqrt((xe-P[1])^2+(ye-P[2])^2)+sqrt((xe-Q[1])^2+(ye-Q[2])^2); min(jum), max(jum)
\end{eulerprompt}
\begin{euleroutput}
  7.54998577103
  8.42613486932
\end{euleroutput}
\begin{eulercomment}
Hasil yang diharapkan: pusat M0 = [3,2.5], c = 2.5, a = 4, b = 3.1225; min dan max jum keduanya 8.
Penjelasan: setiap titik pada kurva memiliki jumlah jarak ke P dan Q yang sama (8), sehingga kurva
tersebut elips dengan fokus P dan Q.

\end{eulercomment}
\eulersubheading{Latihan 5: Hiperbola jika diketahui kedua titik fokus P dan Q}
\begin{eulercomment}
Catatan: pada soal tertulis "ellips ... selisih jarak". Yang dimaksud adalah hiperbola: tempat
kedudukan titik-titik yang SELISIH jarak ke P dan Q-nya (nilai mutlak) konstan.
Uraian materi: hiperbola dengan fokus P, Q: \textbar{}TP - TQ\textbar{} = 2a (0 \textless{} 2a \textless{} PQ). Dengan c = PQ/2 dan
b\textasciicircum{}2 = c\textasciicircum{}2 - a\textasciicircum{}2, pada posisi baku persamaannya x\textasciicircum{}2/a\textasciicircum{}2 - y\textasciicircum{}2/b\textasciicircum{}2 = 1, dengan asimtot y = +/- (b/a) x.

Fokus P(-3,0), Q(3,0), selisih jarak 4 (a=2, c=3, b\textasciicircum{}2=5). Cabang kanan: TP - TQ = 4, cabang kiri: TP - TQ = -4.
\end{eulercomment}
\begin{eulerprompt}
>hipR &= sqrt((x+3)^2+y^2)-sqrt((x-3)^2+y^2)-4; $hipR='0
>hipL &= sqrt((x+3)^2+y^2)-sqrt((x-3)^2+y^2)+4; $hipL='0
>P=[-3,0]; Q=[3,0];
>setPlotRange(8); aspect(1); plotPoint(P,"P"); plotPoint(Q,"Q");
>plot2d(hipR,level=0,add=1,contourcolor=red); plot2d(hipL,level=0,add=1,contourcolor=red):
\end{eulerprompt}
\begin{eulercomment}
Asimtot y = +/- sqrt(5)/2 x:
\end{eulercomment}
\begin{eulerprompt}
>plot2d("sqrt(5)/2*x",add=1); plot2d("-sqrt(5)/2*x",add=1):
\end{eulerprompt}
\begin{eulercomment}
Pemeriksaan numerik dengan titik parametrik (2 cosh t, sqrt(5) sinh t):
\end{eulercomment}
\begin{eulerprompt}
>tt=linspace(-2,2,100); xh=2*cosh(tt); yh=sqrt(5)*sinh(tt);
>sel=abs(sqrt((xh+3)^2+yh^2)-sqrt((xh-3)^2+yh^2)); min(sel), max(sel)
\end{eulerprompt}
\begin{euleroutput}
  4
  4
\end{euleroutput}
\begin{eulercomment}
Hasil yang diharapkan: min dan max keduanya 4.
Penjelasan: setiap titik (2 cosh t, sqrt(5) sinh t) memenuhi x\textasciicircum{}2/4 - y\textasciicircum{}2/5 = 1 dan
selisih jaraknya ke kedua fokus tepat 4, jadi kurva itu hiperbola dengan fokus P dan Q.

\begin{eulercomment}
\eulerheading{Daftar Pustaka dan Catatan}
\begin{eulercomment}
- R. Grothmann, Euler Math Toolbox, dokumentasi dan contoh (file geometry.e, spherical.e, povray.e).
- N. J. Wildberger, Divine Proportions: Rational Trigonometry to Universal Geometry (2005).
- H. S. M. Coxeter dan S. L. Greitzer, Geometry Revisited (1967): garis Euler, titik Fermat,
\end{eulercomment}
\begin{eulerttcomment}
  teorema Pitot, lingkaran dalam dan lingkaran luar.
\end{eulerttcomment}
\begin{eulercomment}

Catatan kejujuran: notebook ini disusun tanpa menjalankannya di EMT. Seluruh nilai hasil yang
diharapkan sudah dihitung dan diperiksa dengan Python. Jika ada perintah yang menghasilkan galat
(misalnya karena perbedaan versi EMT), kemungkinan besar itu perbedaan sintaks kecil; periksa
pesan galatnya dan sesuaikan.
\end{eulercomment}
\end{eulernotebook}
\end{document}
